\chapter{Enhanced traces and sewing}
\label{v4-arith:traces:chapter}

A matrix trace annihilates commutators. The trace of a category can retain
an entire complex and a circle action. An arithmetic comparison must specify
which of these objects it uses, which class it evaluates, and how the
evaluation reaches a scalar partition function. Sewing requires compatible
gluing maps as well.

Two elementary examples distinguish these data. In a matrix algebra, the
identity need not generate degree-zero Hochschild homology. In a Frobenius
algebra, changing the functional can change the sphere value without
changing multiplication. Both phenomena matter when a local vacuum or a
distinguished arithmetic class is part of the comparison.

We begin with ordinary algebras over a field. Tensor products in the first
constructions are algebraic. The later categorical, completed, and Hilbert
constructions will have their own coefficients and domains. The centers and
actions of Chapter~\ref{v4-arith:centers:chapter} act on specified objects.
Passing to their traces is a further operation.

\section{A trace and the class it evaluates}

Let $k$ be a field and $A$ a nonzero unital associative $k$-algebra. Write
$[A,A]$ for the $k$-linear span of commutators $ab-ba$. Define
\[
 HH_0(A)=A/[A,A].
\]
This quotient is a vector space. In general $[A,A]$ is not an ideal, so the
quotient has no multiplication induced by $A$.

A \emph{cyclic functional} on $A$ is a $k$-linear map $\tau:A\to k$
satisfying $\tau(ab)=\tau(ba)$ for every $a,b$. It is equivalent to a linear
functional on $HH_0(A)$. The class $[1_A]$ and the value $\tau(1_A)$ are
separate data. A normalized cyclic functional has $\tau(1_A)=1$.

\begin{proposition}[Matrix trace and the regular object]
\label{v4-arith:traces:matrix-class}
For every integer $d\geq1$, matrix trace induces an isomorphism
\[
 HH_0(M_d(A))\longrightarrow HH_0(A),\qquad
 [X]\longmapsto\left[\sum_i X_{ii}\right].
\]
Its inverse sends $[a]$ to $[aE_{11}]$. Under this isomorphism,
\[
 [1_{M_d(A)}]\longmapsto d[1_A].
\]
Here $E_{ij}$ is the matrix with entry one in position $(i,j)$ and zero
elsewhere.
\end{proposition}
\begin{proof}
The trace of a matrix commutator is a sum of commutators in $A$:
\[
 \operatorname{Tr}(XY-YX)=\sum_{i,j}[X_{ij},Y_{ji}].
\]
Thus trace descends to the quotient.

Conversely, $[a,b]E_{11}=[aE_{11},bE_{11}]$, so the proposed inverse is well
defined. For $i\ne j$,
\[
 aE_{ij}=[aE_{ij},E_{jj}],
\]
and for $i\ne1$,
\[
 aE_{ii}-aE_{11}=[aE_{i1},E_{1i}].
\]
Every matrix class therefore equals the class of its trace placed in the
$(1,1)$ entry. The two maps are inverse. The formula for the identity
follows by adding its $d$ diagonal entries.
\end{proof}

For $A=k$ of characteristic two and $d=2$, the space $HH_0(M_2(k))$ is one
dimensional, but its identity class is zero. Consequently no cyclic
functional on this algebra has value one on the identity. Over a field in
which $d$ is invertible, the normalized matrix functional
$d^{-1}\operatorname{Tr}$ does have that value. The division by $d$ is part
of its normalization.

The matrix comparison has an object-theoretic interpretation. The usual
equivalence between perfect $M_d(A)$-modules and perfect $A$-modules carries
the regular matrix module to $A^{\oplus d}$. Its identity class accordingly
becomes $d[1_A]$. An equivalence of categories does not require these two
regular objects to correspond. An additional object marking can impose that
requirement.

\section{Chains, cyclic rotation, and completion}

Put $\overline A=A/k1$. The normalized Hochschild chains are
\[
 C_n(A)=A\otimes_k\overline A^{\otimes n}\quad(n\geq0),
 \qquad C_n(A)=0\quad(n<0).
\]
Write an elementary tensor as $a_0[a_1|\cdots|a_n]$. An entry equal to $1$
inside the brackets is zero. The differential lowers degree and is
\begin{align}
\label{v4-arith:traces:boundary}
 b\bigl(a_0[a_1|\cdots|a_n]\bigr)
 ={}&a_0a_1[a_2|\cdots|a_n]\\
 &+\sum_{i=1}^{n-1}(-1)^i
     a_0[a_1|\cdots|a_ia_{i+1}|\cdots|a_n]\\
 &+(-1)^na_na_0[a_1|\cdots|a_{n-1}].\notag
\end{align}

The formula in positive degree descends from the unnormalized chains
$A^{\otimes(n+1)}$ by quotienting the tensors with a unit in a bracket
position. Set $b=0$ on degree zero.

\begin{definition}[The regular Hochschild trace]
\label{v4-arith:traces:regular-complex}
The \emph{Hochschild chain complex with regular coefficients} is
$(C_\bullet(A),b)$. Its homology is $HH_\bullet(A)$. The coefficient
bimodule is $A$ itself, with left and right multiplication. Replacing that
bimodule changes the trace complex.
\end{definition}

Associativity proves $b^2=0$. In the unnormalized complex, composing two
multiplication faces in either order gives the same result after renumbering
the later face. Their alternating signs are opposite. This includes the face
multiplying the last entry into the first. In degree zero the resulting
homology is exactly $A/[A,A]$.

The circle operation is already visible on these chains. Define Connes'
operator $B:C_n(A)\to C_{n+1}(A)$ by
\begin{equation}
\label{v4-arith:traces:connes}
 B\bigl(a_0[a_1|\cdots|a_n]\bigr)
 =\sum_{i=0}^{n}(-1)^{ni}
   1[a_i|\cdots|a_n|a_0|\cdots|a_{i-1}].
\end{equation}
For $i=0$, the bracket is $[a_0|a_1|\cdots|a_n]$. Thus
\[
 B(a)=1[a],\qquad
 B(a_0[a_1])=1[a_0|a_1]-1[a_1|a_0].
\]

\begin{proposition}[The mixed complex]
\label{v4-arith:traces:mixed-complex}
The operators satisfy
\[
 b^2=0,\qquad B^2=0,\qquad bB+Bb=0.
\]
In particular the pair of differentials $(b,B)$ contains more information
than $HH_\bullet(A)$ as a graded vector space.
\end{proposition}
\begin{proof}
The multiplication-face cancellation proves the first identity. Every term
of $B$ has leading coefficient one. Applying $B$ again places that
coefficient inside the bracket, so every resulting term is degenerate. This
proves $B^2=0$.

For the remaining identity, expand the boundary of each cyclic rotation in
\eqref{v4-arith:traces:connes}. Its two outer faces cancel the corresponding
outer faces from the adjacent rotation. Each interior multiplication then
occurs once in $bB$ and once in $Bb$. Moving the multiplication past the
cyclic starting point changes its face index by one, with the rotation sign
$(-1)^{ni}$ accounting for the wrap. The signs of these two terms are
opposite. All terms cancel, proving $bB=-Bb$.
\end{proof}

More generally, a \emph{mixed complex} over $k$ is a homologically graded
vector space $C$ with two operators. They have degrees $|b|=-1$, $|B|=1$ and
satisfy these three identities. This definition also applies to the cyclic
complexes of differential graded categories. Their internal differentials
and permutations carry the Koszul signs of their graded morphisms.

For a mixed complex, let $u$ have degree $-2$. Define its negative cyclic
complex by
\begin{equation}
\label{v4-arith:traces:negative-cyclic}
 CC^-_m(C)=\prod_{j\geq0} C_{m+2j}u^j,
 \qquad d=b+uB.
\end{equation}
The mixed identities give $d^2=0$. The product in this formula is part of
the definition.

Its periodic cyclic complex is $CC^{\mathrm{per}}(C)=CC^-(C)[u^{-1}]$, with
the same differential. Inverting $u$ is a filtered colimit after forming the
displayed product. These two limiting operations must be respected by any
comparison involving these complexes.

\begin{example}[A product created by negative cyclic completion]
\label{v4-arith:traces:sum-product}
For $i\geq0$, let $C^{(i)}$ be one copy of $k$ in homological degree $2i$,
with $b=B=0$. Then
\[
 CC^-_0\left(\bigoplus_{i\geq0}C^{(i)}\right)
   =\prod_{i\geq0}k,
 \qquad
 \left(\bigoplus_{i\geq0}CC^-(C^{(i)})\right)_0
   =\bigoplus_{i\geq0}k.
\]
The natural map is the inclusion of finite-support sequences and is not
surjective. All differentials vanish, so the discrepancy also occurs in
homology. An equivariant direct-sum decomposition of Hochschild chains
therefore does not, by itself, give a direct-sum decomposition after
negative cyclic completion.
\end{example}

A cyclic functional $\tau:A\to k$ is a map of mixed complexes from
$C_\bullet(A)$ to $k$ in degree zero, with zero $B$ on the target. The chain
condition is $\tau b(a_0[a_1])=0$. All other degrees give zero, and the
compatibility with $B$ follows from the degrees. This functional is an
evaluation of the mixed complex, rather than the complex itself.

\subsection{Categorical trace and Morita comparison}

For a dg algebra $A$ over a commutative ring $k$, the derived regular trace
is $A\otimes_{A^e}^{\mathrm L}A$. Choose a dg model with $K$-flat morphism
complexes over $k$: tensoring with each complex must preserve
quasi-isomorphisms. A cofibrant replacement provides such a model for the
bar tensors.

This is the Hochschild trace of the identity of its perfect-module category.
The complex, with its circle action, is distinct from a single class
associated to an object. For example a perfect object has its Hochschild
Chern character, giving the additive map $K_0(\operatorname{Perf}A)\to
HH_0(A)$. For the regular object $A$, this class is $[1_A]$.

The derived mixed complex is invariant under dg Morita equivalence. More
precisely, Keller's Theorem 5.2 in \emph{On differential graded categories}
makes the mixed-complex construction a functor on the Morita category and
identifies $C(A)$ with $C(\operatorname{Perf}_{\mathrm{dg}}A)$. This theorem
concerns the enhanced dg categories and their derived equivalences. An
equivalence of lists of simple objects does not meet its hypotheses.

For matrix algebras the comparison has an explicit cyclic-chain formula. On
unnormalized degree-$n$ chains, set
\[
 \operatorname{tr}_{\mathrm{cyc}}
       (X_0\otimes\cdots\otimes X_n)
   =\sum_{i_0,\ldots,i_n}
       (X_0)_{i_0i_1}\otimes
       (X_1)_{i_1i_2}\otimes\cdots\otimes
       (X_n)_{i_ni_0}.
\]

Matrix multiplication proves compatibility with every face, including the
cyclic last face. Inserting an identity matrix collapses an adjacent index
and gives the scalar-unit degeneracy. Rotating the indices proves
compatibility with cyclic rotation. Thus it descends to a map of mixed
complexes. In degree zero it is precisely
Proposition~\ref{v4-arith:traces:matrix-class}.

The corner inclusion used there gives an inverse on $HH_0$. It is not a
unital cyclic inverse on normalized chains. It sends $1_A$ to $E_{11}$,
whereas the matrix-algebra unit is $\sum_iE_{ii}$. The enhanced Morita
comparison accounts for the full units and their homotopies. It preserves
the derived trace while transporting each chosen object marking to the
object actually corresponding under the equivalence.

\section{A functional that supplies sewing}

To contract two boundary entries, a functional needs a nondegenerate
pairing. This condition can be imposed and calculated in finite dimension
before asking for an arithmetic or analytic realization.

\begin{definition}[Commutative Frobenius algebra]
\label{v4-arith:traces:frobenius}
A \emph{commutative Frobenius algebra} over $k$ is a finite-dimensional
commutative unital $k$-algebra $A$ with a linear functional
$\varepsilon:A\to k$ such that
\[
 (a,b)_\varepsilon=\varepsilon(ab)
\]
is nondegenerate. Choose dual bases $e_i,f_i$ for this pairing and put
\[
 c=\sum_i e_i\otimes f_i,\qquad
 E=\sum_i e_if_i.
\]
The tensor $c$ is the inverse pairing, and $E$ is the \emph{handle element}.
Both are independent of the chosen dual bases.
\end{definition}

The multiplication $m$ joins two inputs. The pairing contracts two entries,
and $c$ supplies the reverse contraction. Define
\[
 \Delta(a)=(a\otimes1)c.
\]
This is the linear map adjoint to multiplication for the Frobenius pairing.
In particular its normalization is determined by $\varepsilon$, not by
multiplication alone.

\begin{proposition}[Finite sewing identities]
\label{v4-arith:traces:finite-sewing}
The maps $m$, $\Delta$, the algebra unit and $\varepsilon$ satisfy the unit
and counit identities, associativity and coassociativity, and
\[
 \Delta(ab)=(a\otimes1)\Delta(b)
            =(1\otimes b)\Delta(a).
\]
Furthermore,
\[
 m\Delta(a)=aE.
\]
The closed contraction formed from one unit, $g$ successive handles and one
counit is
\[
 Z_g=\varepsilon(E^g)\qquad(g\geq0).
\]
\end{proposition}
\begin{proof}
The inverse-pairing identity is $\sum_i\varepsilon(af_i)e_i=a$. It proves
the counit identity and shows that the adjoint of $m$ is the displayed
$\Delta$. Taking adjoints of associativity proves coassociativity, since the
tensor-product pairing is nondegenerate.

The inverse tensor obeys $(a\otimes1)c=(1\otimes a)c$. To verify this
equality, pair both tensors with $x\otimes y$. Each value is
$\varepsilon(axy)$. Nondegeneracy proves the equality, and then proves both
Frobenius identities for $\Delta(ab)$. Multiplying the two factors of
$\Delta(a)$ gives $aE$. Iterating this handle map on $1$ and applying
$\varepsilon$ gives the formula for $Z_g$.
\end{proof}

\begin{example}[Semisimple weights]
\label{v4-arith:traces:semisimple-weights}
Let $A=k^r$, with orthogonal primitive idempotents $e_i$. Choose nonzero
$w_i\in k$ and set $\varepsilon(e_i)=w_i$. The dual basis is $w_i^{-1}e_i$,
so
\[
 \Delta(e_i)=w_i^{-1}e_i\otimes e_i,\qquad
 E=\sum_i w_i^{-1}e_i,\qquad
 Z_g=\sum_i w_i^{1-g}.
\]
In particular $Z_0=\sum_i w_i$ and $Z_1=r$ in $k$. The algebra does not
determine the weights. The counit and the inverse pairing retain them
through every contraction.
\end{example}

\begin{example}[A nilpotent algebra with a nondegenerate pairing]
\label{v4-arith:traces:nilpotent-frobenius}
Let $A=k[\epsilon]/(\epsilon^2)$ and choose $\varepsilon(1)=a$,
$\varepsilon(\epsilon)=1$. Its Gram matrix in the basis $(1,\epsilon)$ is
\[
 \begin{pmatrix}a&1\\1&0\end{pmatrix},
\]
whose determinant is $-1$ in every characteristic. The dual basis is
$(\epsilon,1-a\epsilon)$. Consequently
\[
 E=2\epsilon,\qquad Z_0=a,\qquad Z_1=2,
 \qquad Z_g=0\quad(g\geq2).
\]

These are algebraic contractions. Over $\mathbb R$, the displayed Gram
matrix is indefinite. Nondegeneracy supplies sewing, but does not assert
positivity of a Hilbert-space pairing.
\end{example}

The finite identities specify what must be transported in a larger
comparison: the functional, inverse pairing, unit, and gluing maps. An
equality of closed scalar values cannot recover all these maps. For
infinite-dimensional or completed objects, defining the inverse pairing and
proving convergence of its contractions are additional parts of the
construction.

\section{Spherical traces and marked closures}

A spherical category supplies a trace by its duality maps. It need not
supply a cyclic functional on an algebra used to present its objects. The
arithmetic fusion example makes this distinction explicit.

Let $\mathcal C$ be a split spherical fusion category over a
characteristic-zero field $k$. Thus $\mathcal C$ is a rigid semisimple
$k$-linear tensor category with finitely many simple objects,
finite-dimensional morphism spaces, and simple unit. Splitness means
$\operatorname{End}(S)=k$ for each simple $S$. Its chosen pivotal duality
identifies left and right categorical traces. Write $S_i$ for its simple
representatives and
\[
 d_i=\operatorname{tr}_{S_i}(1_{S_i})
\]
for their spherical dimensions. For an object $X\simeq\bigoplus_i
S_i^{\oplus m_i}$, an endomorphism has multiplicity matrices $f_i\in
M_{m_i}(k)$.

\begin{proposition}[Trace on multiplicity spaces]
\label{v4-arith:traces:spherical-matrices}
The categorical trace is
\[
 \operatorname{tr}_X(f)=\sum_i d_i\operatorname{Tr}_k(f_i).
\]
In particular it is a cyclic functional on $\operatorname{End}(X)$. The
coefficients $d_i$ are determined by the chosen duality.
\end{proposition}
\begin{proof}
Choose inclusions $u_{i\alpha}:S_i\to X$ and projections $v_{i\alpha}:X\to
S_i$ with
\[
 v_{i\alpha}u_{j\beta}
   =\delta_{ij}\delta_{\alpha\beta}1_{S_i},
 \qquad
 \sum_{i,\alpha}u_{i\alpha}v_{i\alpha}=1_X.
\]
Expand $f$ using these splittings. Cyclicity of the spherical trace gives
$\operatorname{tr}_X(u_{i\alpha}v_{i\beta})
=\operatorname{tr}_{S_i}(v_{i\beta}u_{i\alpha}) =\delta_{\alpha\beta}d_i$.
The diagonal matrix entries therefore give the displayed sum. Ordinary
matrix trace is cyclic, proving the last assertion.
\end{proof}

For the order-two arithmetic category of Chapter~\ref{v4-arith:mf:chapter},
the fusion rules are
\[
 g^2=1,\qquad g\rho=\rho g=\rho,\qquad
 \rho^2=1+g+2\rho.
\]
The constructed spherical structures have $d_1=d_g=1$ and
\[
 d_\rho=1-\sqrt3
 \quad\hbox{or}\quad
 d_\rho=1+\sqrt3
\]
at the two indicated embeddings of the coefficient field. Indeed
multiplicativity and additivity of dimension require $d_\rho^2=2+2d_\rho$.
Their global dimensions and traces on the sum of the simples are
respectively
\[
 \sum_i d_i^2=6\mp2\sqrt3,\qquad
 \operatorname{tr}_{1\oplus g\oplus\rho}(1)=3\mp\sqrt3.
\]

These are different contractions. The unitary realization in
Chapter~\ref{v4-arith:centers:chapter} uses $1+\sqrt3$. The original
arithmetic embedding retains $1-\sqrt3$ and the nonunitary realization of
Chapter~\ref{v4-arith:asi:chapter}.

One cannot change the sign of a single dimension while keeping that category
and its pivotal structure. More precisely, any tensor automorphism of the
identity functor acts on the three simples by scalars
$\eta_1=1,\eta_g,\eta_\rho$. The unit and $\rho$ summands of $\rho^2$
require
\[
 \eta_\rho^2=1,\qquad \eta_\rho^2=\eta_\rho.
\]
Hence $\eta_\rho=1$, and the $g$ summand gives $\eta_g=1$. Two pivotal
structures differ by such an automorphism, so the existing pivotal structure
is unique. Passing to the other field embedding changes the coefficient data
of the category.

The algebra $\operatorname{End}(\bigoplus_i S_i)$ is $k^r$. When every $d_i$
is nonzero, its spherical trace makes this commutative algebra a Frobenius
algebra with weights $w_i=d_i$.
Example~\ref{v4-arith:traces:semisimple-weights} then gives its finite
contractions. This construction records a trace on the sum of the simples.
Recovering the fusion tensor, its associator or its geometric sewing
requires those maps as additional data.

\subsection{The scalar unit of a Leavitt algebra}

Let $L_r$ be the unital algebra with generators
$s_1,\ldots,s_r,t_1,\ldots,t_r$ and relations
\[
 t_is_j=\delta_{ij}1,\qquad \sum_i s_it_i=1.
\]
The $t_i$ are algebraic dual generators. No complex conjugation is part of
this notation.

\begin{proposition}[A vanishing unit class]
\label{v4-arith:traces:leavitt-unit}
If $r-1$ is invertible in $k$, then $[1]=0$ in $HH_0(L_r)$. Consequently
every cyclic functional on $L_r$ vanishes at $1$.
\end{proposition}
\begin{proof}
The defining relations give the literal commutator identity
\[
 \sum_{i=1}^r[t_i,s_i]
   =r1-\sum_i s_it_i=(r-1)1.
\]
Divide by $r-1$ and pass to $L_r/[L_r,L_r]$.
\end{proof}

This argument concerns one class. It does not assert that $HH_0(L_r)$ is
zero. It also has the stated coefficient restriction. If $r=1$, the algebra
is $k[s,s^{-1}]$, and the identity class is nonzero.

For the arithmetic algebra $\mathcal B_I$ of
Section~\ref{v4-arith:mf:leavitt}, one has $r=2N$. Its characteristic-zero
coefficient field makes $2N-1$ invertible. Nevertheless its marked
categorical closure recovers the arithmetic values. Retain the object
$\rho$, its duality maps $X_0,X'_0$, and the scalar network $\mathcal
N_g=b_g1_\rho$. With $d=-1/b_0$ and $e_g=b_g/b_0$, the closure is
\[
 \operatorname{cl}_\rho(a)=dX'_0\rho(a)X_0,\qquad
 -\operatorname{cl}_\rho(\mathcal N_g)=e_g.
\]

The second identity follows from $X'_0X_0=1$ and the fact that $\rho$ fixes
scalars. This functional has domain $\operatorname{End}(\rho)$ and uses its
specified duality. Its value is compatible with the vanishing ambient unit
class because it is a different trace construction.

The finite balanced-tensor realization of
Theorem~\ref{v4-arith:mf:finite-functor} retains this network as well. Its
matrix detector is
\[
 M_g=b_g1_{\omega(\rho)},\qquad
 b_g=\frac{\operatorname{Tr}_k(M_g)}{3N}.
\]
This formula uses the scalar endomorphism and the particular dimension
$\dim_k\omega(\rho)=3N$. It is not a replacement for categorical trace on
arbitrary diagrams. The functor takes values in bimodules over $k^{N+1}$,
with their balanced tensor. It is not a finite-dimensional unital
representation of $\mathcal B_I$.

\section{Extensions and regular Hochschild traces}

The trace of a nonsemisimple category requires its actual complex.
Extensions between simple modules give useful information, but they need not
survive regular Hochschild trace. Two small algebras show both behaviors.

Let $A$ be a finite-dimensional algebra over a field $k$. There are finitely
many simple right $A$-modules up to isomorphism: each occurs among the
composition factors of the regular module. If $M,N$ are finite-dimensional,
then $\operatorname{Ext}^n_A(M,N)$ is finite-dimensional for each $n$.
Indeed successive surjections from finite free modules give a projective
resolution with finite-dimensional terms. Applying
$\operatorname{Hom}_A(-,N)$ proves degreewise finiteness. There is no
resulting bound on the total graded dimension.

\subsection{The dual numbers}

Put $A=k[\epsilon]/(\epsilon^2)$ and $S=A/(\epsilon)$. Multiplication by
$\epsilon$ has kernel and image $k\epsilon$, so
\[
 \cdots\xrightarrow{\epsilon} A\xrightarrow{\epsilon} A
 \xrightarrow{\epsilon} A\longrightarrow S\longrightarrow0
\]
is an exact free resolution. Applying $\operatorname{Hom}_A(-,S)$ makes
every differential zero. Thus
\[
 \operatorname{Ext}^n_A(S,S)=k\qquad(n\geq0).
\]
The algebra has radical square zero, but its simple module has infinite
projective dimension. In characteristic two it is also $k[C_2]$, under
$\epsilon\mapsto g-1$.

For the regular trace, use the enveloping algebra $A^e=A\otimes_k
A^{\mathrm{op}}$ and the coefficient bimodule $A$. The bar resolution
identifies
\[
 HH_n(A)=\operatorname{Tor}^{A^e}_n(A,A).
\]
Here both occurrences of $A$ retain their regular actions.

\begin{proposition}[The complete regular trace of the dual numbers]
\label{v4-arith:traces:dual-numbers}
For every field $k$,
\[
 HH_0(A)=A,\qquad
 HH_{2r+1}(A)=A/(2\epsilon A)\quad(r\geq0),\qquad
 HH_{2r}(A)=\operatorname{ann}_A(2\epsilon)\quad(r\geq1).
\]
If $\operatorname{char}k\ne2$, each positive-degree group has dimension one.
If $\operatorname{char}k=2$, then $HH_n(A)=A$ in every nonnegative degree.
\end{proposition}
\begin{proof}
Write $A^e=k[x,y]/(x^2,y^2)$, with $x$ acting on the left and $y$ on the
right. Set $u=x-y$ and $v=x+y$. There is a free bimodule complex
\[
 \cdots\xrightarrow{u}A^e\xrightarrow{v}A^e
 \xrightarrow{u}A^e\xrightarrow{\mu}A\longrightarrow0,
 \qquad \mu(x)=\mu(y)=\epsilon.
\]
The differential out of odd degree is $u$ and that out of positive even
degree is $v$. Their composites vanish.

For $z=\alpha+\beta x+\gamma y+\eta xy$,
\[
 uz=\alpha(x-y)+(\gamma-\beta)xy,\qquad
 vz=\alpha(x+y)+(\beta+\gamma)xy.
\]
Consequently
\[
 \ker u=\langle x+y,xy\rangle=\operatorname{im}v,\qquad
 \ker v=\langle x-y,xy\rangle=\operatorname{im}u.
\]

Also $\ker\mu=\langle x-y,xy\rangle$. These statements remain valid in
characteristic two, and prove exactness.

Tensoring with the regular $A$ sends $u$ to zero and $v$ to $2\epsilon$. The
resulting alternating differential gives the asserted quotients and
annihilators. If two is invertible, the annihilator is $\epsilon A$. If two
is zero, every differential vanishes.

The normalized Hochschild formula gives the same calculation without a
change of resolution. Since $\overline A=k\epsilon$, interior products in a
normalized chain vanish, and the two outer faces give
\[
 b_n\bigl(c[\epsilon|\cdots|\epsilon]\bigr)
  =(1+(-1)^n)c\epsilon[\epsilon|\cdots|\epsilon].
\]
The bracket on the right has $n-1$ entries. In particular the cyclic last
face is responsible for the factor two in even degree.
\end{proof}

In characteristic two, the simple-module Ext groups have dimension one while
the regular Hochschild groups have dimension two. The coefficient module and
the traced resolution are indispensable parts of either statement.

\subsection{A directed algebra and cyclic endpoint matching}

Let $Q$ be a finite quiver, $J$ its arrow ideal, and $B=kQ/J^2$. Concatenate
paths from left to right and use right modules. For a vertex $i$, put
$P_i=e_iB$ and $S_i=e_iB/e_iJ$. The radical of $P_i$ is one copy of $S_j$
for each arrow $i\to j$. Iterating projective covers therefore gives
\begin{equation}
\label{v4-arith:traces:path-ext}
 \dim_k\operatorname{Ext}^n_B(S_i,S_j)
  =\#\{\text{oriented length-}n\text{ paths from }i\text{ to }j\}.
\end{equation}

To justify the last step, every differential in this minimal resolution
lands in a radical. A map to a simple module kills that radical, so the Hom
complex has zero differential. All groups of degree at least two vanish
precisely when the quiver has no composable pair of arrows.

Take
\[
 Q=(0\xrightarrow{a}1\xrightarrow{b}2),\qquad B=kQ/(ab).
\]
Then $P_0=\langle e_0,a\rangle$, $P_1=\langle e_1,b\rangle$, and
$P_2=\langle e_2\rangle$. The sequence
\[
 0\longrightarrow P_2\xrightarrow{e_2\mapsto b}P_1
 \xrightarrow{e_1\mapsto a}P_0\longrightarrow S_0
 \longrightarrow0
\]
is exact: the middle kernel is $kb$ because $ab=0$. It gives
\[
 \operatorname{Ext}^1_B(S_0,S_1)=k,\qquad
 \operatorname{Ext}^1_B(S_1,S_2)=k,\qquad
 \operatorname{Ext}^2_B(S_0,S_2)=k.
\]

The product of the two adjacent degree-one classes is nonzero. The exact
sequence $0\to S_1\to P_0\to S_0\to0$ identifies its Yoneda product with the
dimension-shifting isomorphism $\operatorname{Ext}^1_B(S_1,S_2)\simeq
\operatorname{Ext}^2_B(S_0,S_2)$. Thus the algebra relation $ab=0$ does not
make this product zero.

For comparison, the path algebra $kQ$ before imposing $ab=0$ has
$\operatorname{rad}P_0\simeq P_1$ and $\operatorname{rad}P_1\simeq P_2$. Its
simple modules have projective resolutions of length at most one. It is a
different algebra, and its $P_0$ has Loewy length three.

\begin{proposition}[Regular trace of the directed model]
\label{v4-arith:traces:directed-model}
For $B=kQ/(ab)$ as above,
\[
 HH_0(B)=ke_0\oplus ke_1\oplus ke_2,\qquad
 HH_n(B)=0\quad(n\geq1).
\]
\end{proposition}
\begin{proof}
Set $S=ke_0\oplus ke_1\oplus ke_2$ and $J=ka\oplus kb$. The relative
normalized bar terms $P_n=B\otimes_S J^{\otimes_S n}\otimes_S B$ are
\begin{align*}
 P_0&=\bigoplus_{i=0}^2 Be_i\otimes_k e_iB,\\
 P_1&=(Be_0\otimes_k e_1B)\oplus(Be_1\otimes_k e_2B),\\
 P_2&=Be_0\otimes_k e_2B,
\end{align*}
and $P_n=0$ for $n\geq3$. Each summand is a projective bimodule, being an
idempotent summand of $B\otimes_k B$. The differentials are
\begin{align*}
 d_1(l\otimes j\otimes r)&=lj\otimes r-l\otimes jr,\\
 d_2(l\otimes a\otimes b\otimes r)
   &=la\otimes b\otimes r+l\otimes a\otimes br.
\end{align*}
The interior term of $d_2$ vanishes because $ab=0$.

For completeness, the dimensions of these terms are $8,4,1$. A basis of
$P_1$ consists of
\begin{align*}
 q_1&=e_0\otimes a\otimes e_1,&
 q_2&=e_0\otimes a\otimes b,\\
 q_3&=e_1\otimes b\otimes e_2,&
 q_4&=a\otimes b\otimes e_2.
\end{align*}
Their images are
\[
 a\otimes e_1-e_0\otimes a,\quad a\otimes b,\quad
 b\otimes e_2-e_1\otimes b,\quad -a\otimes b.
\]
The first three are independent. The kernel of $d_1$ is $k(q_2+q_4)$, and
$d_2$ sends its generator to $q_2+q_4$.

Multiplication $P_0\to B$ is surjective with a three-dimensional kernel.
This proves exactness in every characteristic.

Regular trace sends each summand to
\[
 B\otimes_{B^e}(Be_i\otimes_k e_jB)\simeq e_jBe_i.
\]
The forward map is $c\otimes(l\otimes r)\mapsto rcl$. The inverse is
$t\mapsto t\otimes(e_i\otimes e_j)$. The balancing relations prove that
these maps are inverse.

All off-diagonal terms here vanish, because a directed path has no return
path. The diagonal terms are $ke_i$. The traced complex is consequently
$k^3$ in degree zero and zero in every positive degree.
\end{proof}

The relation contributes a nonzero degree-two bimodule-resolution term.
Taking the regular trace kills that term by its endpoints. An arithmetic use
of simple-module extensions must therefore construct the map into its trace
complex and determine the image of the Yoneda products.

\section{Tensoring local trace complexes}

Local degree bounds do not simply add their positive parts. Fix a
commutative ring $R$. We use homological shifts, so a module $M[1]$ is
concentrated in degree one. Suppose two complexes have specified splittings
\[
 C_2\simeq B_2\oplus W_2[1],\qquad
 C_3\simeq B_3\oplus W_3[1],
\]
Here $B_2,W_2,B_3,W_3$ are $R$-modules in degree zero. The subscripts label
two local factors.

\begin{proposition}[All four tensor terms]
\label{v4-arith:traces:four-terms}
There is an equivalence
\begin{align*}
 C_2\otimes_R^{\mathrm L}C_3\simeq{}&
 B_2\otimes_R^{\mathrm L}B_3\\
 &\oplus(W_2\otimes_R^{\mathrm L}B_3)[1]\\
 &\oplus(B_2\otimes_R^{\mathrm L}W_3)[1]\\
 &\oplus(W_2\otimes_R^{\mathrm L}W_3)[2].
\end{align*}
If all remaining local factors assemble into a degree-zero complex $D$,
their contribution is the derived tensor of $D$ with every displayed term.
\end{proposition}
\begin{proof}
Choose projective resolutions of the modules. Distribute the tensor product
over the two finite direct sums. The tensor differential is $d(x\otimes
y)=dx\otimes y+(-1)^{|x|}x\otimes dy$, and the two shifts in the last term
add to two. Associativity of derived tensor gives the assertion about $D$.
\end{proof}

Over a field, take $B_2=B_3=k$, $W_2=k$, and $W_3=k^2$. The result is
\[
 k\oplus k^3[1]\oplus k^2[2].
\]
If $x$ spans $W_2$ and $y,z$ form a basis of $W_3$, then $x\otimes y$ and
$x\otimes z$ are independent degree-two cycles. Every differential is zero,
so they are not boundaries. Anticommutativity does not cancel products in
distinct factors. Taking unit reference vectors at almost all primes still
permits non-unit entries at both selected primes.

Integral coefficients can give a different result, for a specific algebraic
reason.

\begin{proposition}[Distinct-primary vanishing]
\label{v4-arith:traces:primary-vanishing}
If $M_2$ is a $2$-primary torsion abelian group and $M_3$ is a $3$-primary
torsion abelian group, then
\[
 M_2\otimes_{\mathbb Z}^{\mathrm L}M_3=0.
\]
If $B_2,B_3$ are flat abelian groups, the tensor of $B_2\oplus M_2[1]$ and
$B_3\oplus M_3[1]$ is concentrated in degrees zero and one.
\end{proposition}
\begin{proof}
The free resolution of $\mathbb Z/2^a$ is $[\mathbb
Z\xrightarrow{2^a}\mathbb Z]$. After tensoring with $\mathbb Z/3^b$, its
differential is invertible. Its tensor product and first Tor therefore
vanish. Finite primary groups are sums of such cyclic groups. Every primary
torsion abelian group is the filtered union of its finite subgroups, and
tensor and Tor commute with these filtered colimits.

There is no higher Tor over $\mathbb Z$: the kernel of a surjection from a
free abelian group is free. This proves the first assertion. Flatness
removes higher Tor from the other three terms of
Proposition~\ref{v4-arith:traces:four-terms}.
\end{proof}

For $B_2=B_3=\mathbb Z$, $M_2=\mathbb Z/2$, and $M_3=(\mathbb Z/3)^2$, the
result is
\[
 \mathbb Z\oplus
   \bigl(\mathbb Z/2\oplus(\mathbb Z/3)^2\bigr)[1].
\]
This calculation does not replace the common-field example. They have
different coefficients. An arithmetic comparison needs its integral models,
change-of-coefficient maps, and a proof that the chosen completion preserves
the asserted tensor.

\begin{question}[Arithmetic trace assembly]
\label{v4-arith:traces:arithmetic-assembly}
Construct the local arithmetic categories and regular trace complexes, the
maps from their extension data, and the completed restricted tensor that
assembles them globally. Determine whether the mixed term survives, vanishes
by a coefficient argument, or is removed by a specified quotient. In the
last case, construct the quotient map and prove its compatibility with the
trace operations. A degree-one truncation and the full trace have different
comparison problems.
\end{question}

\section{Orthogonal blocks and cyclic decomposition}

Let $k$ now be a commutative ring, and let $(\mathcal A_i)_{i\in I}$ be
small $k$-linear differential graded categories with $K$-flat morphism
complexes. Here $K$-flat means that tensoring with the complex preserves
quasi-isomorphisms. This hypothesis lets ordinary tensor products compute
the derived tensors used below.

Form the orthogonal union $\mathcal A$: objects belong to one of the
$\mathcal A_i$, and the morphism complex between different blocks is zero.
Let $\operatorname{Perf}(\mathcal A)$ be the category obtained from its
representable modules by finite sums, shifts, cones and retracts.
Write $C_\bullet(\mathcal B)$ for the derived Hochschild mixed complex
of a dg category $\mathcal B$.

\begin{proposition}[Cyclic block decomposition]
\label{v4-arith:traces:blocks}
There is an $S^1$-equivariant equivalence
\[
 C_\bullet(\operatorname{Perf}(\mathcal A))
 \simeq\bigoplus_{i\in I}
       C_\bullet(\operatorname{Perf}(\mathcal A_i)).
\]
Every perfect module has finite block support.
\end{proposition}
\begin{proof}
In cyclic degree $n$, the bar construction is the sum of
\[
 \mathcal A(X_n,X_0)\otimes_k
 \mathcal A(X_{n-1},X_n)\otimes_k\cdots\otimes_k
 \mathcal A(X_0,X_1)
\]
over its cyclic words of objects. A word using two blocks has a zero
morphism factor, including when the change of block occurs at the closing
pair. Every surviving word therefore lies in one block. Composition faces,
unit degeneracies, internal differentials and cyclic rotation preserve it.

Thus the cyclic object itself decomposes as the direct sum of the block
cyclic objects. The $K$-flat assumption computes the derived tensors, and
realization commutes with direct sums. The cyclic structure, hence the
circle action, is preserved. Passing to perfect modules is the dg Morita
invariance of this derived trace.

A representable is supported in one block. Each finite sum, shift, cone or
retract preserves finite support, so every perfect object has it.
Conversely, a finite-support family of perfect block modules is perfect over
$\mathcal A$. This identifies the compact model used in the formula.
\end{proof}

For a compactly generated presentable category, this statement applies when
its compact objects have the specified model. A product description of a
presentable category alone does not identify the relevant trace or its
compact support. Moreover the proposition is an equivalence of cyclic
complexes. A multiplicative trace comparison needs a compatible monoidal
enhancement and preservation of its operations.

Even coordinatewise multiplication on the right has a unit only under an
additional finiteness condition. If infinitely many blocks are nonzero
unital algebras, every element of their direct sum omits some block. It
cannot act as identity on that block. The infinite coordinatewise direct sum
is then nonunital. Finally, negative cyclic completion need not preserve
this sum, as Example~\ref{v4-arith:traces:sum-product} shows.

For arithmetic block comparisons, the coefficient ring, dg models and
indexing maps must be given before applying this decomposition. The
perfect-stack loop comparisons of Chapter~\ref{v4-arith:ic:chapter} also
retain their own geometric hypotheses. The topological circle, a finite
cyclic classifying stack, and a profinite mapping object define different
constructions. Passing between them requires comparison maps and continuity
theorems on the chosen carriers.

\subsection{Filtered continuity on a specified trace carrier}

The same finite-word argument applies to a filtered union of dg categories.
Every cyclic word, morphism tensor and relation appears at a common finite
stage. For $K$-flat models over a common ring, tensoring and filtered
colimits commute. Thus the cyclic complex of that union is the filtered
colimit of its cyclic complexes, before taking homology or cyclic
completion.

There is also a categorical version with a different output. Let $\mathcal
V$ be a presentable symmetric monoidal category whose tensor preserves
colimits separately. For an associative algebra object $A$ of $\mathcal V$,
define
\[
 \operatorname{Tr}_{\mathcal V}(A)
    =\left|[n]\longmapsto A^{\otimes(n+1)}\right|.
\]
The faces multiply adjacent entries, including the cyclic closing pair, and
the degeneracies insert the unit. The bars denote geometric realization of
the cyclic object.

For $\mathcal V=D(k)$ this is a derived Hochschild complex. For $\mathcal
V=\mathrm{Pr}^{\mathrm L}_{k,\mathrm{st}}$ it is a category, as in
Section~\ref{v4-arith:ic:traces}.

\begin{proposition}[Filtered continuity of cyclic realization]
\label{v4-arith:traces:filtered-continuity}
Let $(A_i)_{i\in I}$ be a filtered diagram of associative algebra objects in
$\mathcal V$, and let $A$ be its colimit. Then there is a natural
circle-equivariant equivalence
\[
 \operatorname*{colim}_i\operatorname{Tr}_{\mathcal V}(A_i)
       \simeq\operatorname{Tr}_{\mathcal V}(A).
\]
If the diagram has compatible $E_2$ structures, this equivalence preserves
the induced $E_1$ multiplication on the trace.
\end{proposition}
\begin{proof}
Filtered colimits of associative algebras are computed on their underlying
objects. Indeed all operations have finite arity, and the diagonal $I\to
I^m$ is cofinal for every finite $m$. Separate preservation of colimits
therefore gives
\[
 A^{\otimes(n+1)}
   \simeq\operatorname*{colim}_i A_i^{\otimes(n+1)}.
\]
These equivalences preserve multiplication, the unit and permutations. They
consequently identify every cyclic face, degeneracy and rotation.

Geometric realization is a colimit and commutes with the filtered colimit.
Realizing the degreewise equivalence proves the assertion with its circle
action. For $E_2$ inputs, use their description as associative algebras in
associative algebras. The cyclic construction in one direction retains the
multiplication in the other, with its interchange coherences. The same
colimit argument preserves these operations and proves the last assertion.
\end{proof}

To apply this result to an arithmetic formal presentation, one must identify
its category with the stated colimit of its finite-stage categories. The
transition functors must preserve the convolution structure being used. A
closed immersion of parameter spaces does not by itself establish this
categorical equivalence or its support conditions.

After that input, the trace-continuity step is positive. If $T_i\simeq Q_i$
are compatible finite-stage trace comparisons and
$\operatorname*{colim}_iQ_i\simeq Q$ is independently proved, their colimits
give $T\simeq Q$. The loop-space computation and passage of a
singular-support condition to $Q$ remain separate hypotheses. Functions on
an inverse limit, a pro-\'etale mapping object, and negative cyclic
completion cannot be exchanged with these colimits without their own
arguments.

An arithmetic determinant trace is different again. The trace coefficient of
a universal Chenevier determinant is a function on a deformation ring.
Identifying it with the invariant trace of a universal
$(\varphi,\Gamma)$-module requires a comparison of those coefficient
objects. Compatibility with a specified Nygaard filtration is an additional
equation on that comparison. Neither scalar Hochschild continuity nor the
categorical trace theorem constructs it.

\section{Arithmetic shifts, phases, and occupation traces}

The Bost--Connes algebra gives a concrete arithmetic representation. Its
occupation trace, algebraic Hochschild cochains, and smooth Hochschild
cochains can all be studied. Relating any of them to an independently
defined arithmetic trace category is a further comparison.

\begin{definition}[The Bost--Connes algebra]
\label{v4-arith:traces:bc-algebra}
Let $A_{\mathrm{BC}}$ be the complex unital $C^*$-algebra generated by
isometries $\mu_n$, for $n\geq1$, and unitaries $e(\gamma)$, for
$\gamma\in\mathbb Q/\mathbb Z$, with relations
\begin{align*}
 \mu_1&=1,& \mu_m\mu_n&=\mu_{mn},&
 \mu_n^*\mu_n&=1,\\
 e(0)&=1,& e(\gamma)e(\eta)&=e(\gamma+\eta),&
 e(\gamma)^*&=e(-\gamma),\\
 \mu_m^*\mu_n&=\mu_n\mu_m^* &&\text{if }(m,n)=1,\\
 \mu_ne(\gamma)\mu_n^*
   &=\frac1n\sum_{n\delta=\gamma}e(\delta).
\end{align*}
In particular $\mu_n\mu_n^*=e_n=n^{-1}\sum_{n\delta=0}e(\delta)$. The time
evolution is
\[
 \sigma_t(\mu_n)=n^{it}\mu_n,\qquad
 \sigma_t(e(\gamma))=e(\gamma).
\]
\end{definition}

On $\mathcal L=\ell^2(\mathbb N^\times)$, with orthonormal basis $\delta_m$,
define
\[
 \mu_n\delta_m=\delta_{nm},\qquad
 e(\gamma)\delta_m=\exp(2\pi i m\gamma)\delta_m.
\]
We use the same notation for the abstract generators and these operators
once faithfulness is proved.

\begin{proposition}[The faithful occupation representation]
\label{v4-arith:traces:bc-faithful}
These operators define a faithful representation of $A_{\mathrm{BC}}$ on
$\mathcal L$. The range projection $\mu_n\mu_n^*$ selects the integers
divisible by $n$. Its spectrum is $\{0,1\}$ for $n>1$.
\end{proposition}
\begin{proof}
The adjoint shift sends $\delta_m$ to $\delta_{m/n}$ if $n\mid m$, and to
zero otherwise. This gives the isometry, semigroup and coprime-adjoint
relations. The transfer relation is the finite character identity
\[
 \frac1n\sum_{n\delta=\gamma}\exp(2\pi i m\delta)
  =
  \begin{cases}
   \exp(2\pi i(m/n)\gamma),&n\mid m,\\
   0,&n\nmid m.
  \end{cases}
\]
It also gives the range projection. For $n>1$, its range contains $\delta_n$
and its kernel contains $\delta_1$, so both spectral values occur.

To prove faithfulness, let $K=\prod_p U(1)$, with the product topology. For
$z\in K$, write $z(n)=\prod_p z_p^{v_p(n)}$. The gauge action $\alpha_z$
fixes the phases and sends $\mu_n$ to $z(n)\mu_n$. It preserves the defining
relations.

The finite normal sums
\[
 A_{\mathrm{alg}}
  =\operatorname{span}_{\mathbb C}
     \{\mu_r e(\gamma)\mu_s^*:r,s\geq1,\ %
                              \gamma\in\mathbb Q/\mathbb Z\}
\]
form a dense star subalgebra. Indeed the transfer identities give
$e(\gamma)\mu_n=\mu_ne(n\gamma)$ and its adjoint. If $s=ds'$, $u=du'$, and
$(s',u')=1$, multiplication becomes
\[
 (\mu_re(\gamma)\mu_s^*)(\mu_ue(\eta)\mu_v^*)
   =\mu_{ru'}e(u'\gamma+s'\eta)\mu_{vs'}^*.
\]
Adjoint exchanges $r,s$ and negates $\gamma$. The gauge orbit of each normal
term is continuous. Density and isometry prove continuity on every element
of $A_{\mathrm{BC}}$.

A normal term has gauge character $z(r/s)$. Haar averaging therefore kills
it unless $r=s$, in which case the transfer relation puts it in the phase
algebra $D$. Thus $E(a)=\int_K\alpha_z(a)\,dz$ is a conditional expectation
onto $D$. The representation of $C^*(\mathbb Q/\mathbb Z)=C(\widehat{\mathbb
Z})$ on $\mathcal L$ evaluates at the positive integers. Those integers meet
every residue class and are dense in $\widehat{\mathbb Z}$. Their
evaluations are faithful, so $D=C(\widehat{\mathbb Z})$ and the
representation is faithful on $D$.

The expectation itself is faithful. For $a\geq0$ and a state $\omega$, the
function $z\mapsto\omega(\alpha_z(a))$ is continuous and nonnegative. If
$E(a)=0$, its integral vanishes. Full support of Haar measure makes its
value at the identity zero. States separate positive elements, so $a=0$.

Finally, $V_z\delta_m=z(m)\delta_m$ implements $\alpha_z$ in the occupation
representation $\pi$. If $\pi(a)=0$, then $\pi(E(a^*a))=0$. Faithfulness on
$D$, followed by faithfulness of $E$, gives $a^*a=0$ and $a=0$.
\end{proof}

\subsection{Normalized creation and higher modes}

For each prime $p$, let $\mathscr F_p=\ell^2(\mathbb N_0)$, with number
states $|r\rangle_p$. Form the Hilbert restricted tensor product with vacuum
reference vectors $|0\rangle_p$. On finite occupation vectors, define
\[
 a_p|r\rangle_p=\sqrt r\,|r-1\rangle_p,\qquad
 a_p^\dagger|r\rangle_p=\sqrt{r+1}\,|r+1\rangle_p,\qquad
 N_p=a_p^\dagger a_p.
\]

Unique prime factorization gives a unitary
\[
 U:\mathcal L\longrightarrow
       \widehat{\bigotimes}'_p\mathscr F_p,\qquad
 U\delta_m=\bigotimes_p|v_p(m)\rangle_p.
\]
The arithmetic prime shift is
\[
 U\mu_pU^{-1}=a_p^\dagger(N_p+1)^{-1/2}.
\]
The diagonal factor cancels the creation coefficient, leaving an isometry.
This identity begins on finite occupation vectors and extends boundedly. The
unnormalized creation operator is unbounded.

For $z\in\mathbb C$, the coherent vector
\[
 \exp(-|z|^2/2)\sum_{k\geq0}
     \frac{z^k}{\sqrt{k!}}U\delta_{p^k}
\]
has norm one. It lies in the arithmetic algebra orbit of the vacuum: apply
the norm-convergent series
$\exp(-|z|^2/2)\sum_{k\geq0}z^k(U\mu_pU^{-1})^k/\sqrt{k!}$. The ratio test
proves convergence of its scalar coefficient norms.

For a finite prime set $F$, consider the bosonic current Fock space
\[
 \mathscr F_F^{\mathrm{cur}}
   =\Gamma_s\bigl(\ell^2(F\times\mathbb N_{\geq1})\bigr).
\]
Here $\Gamma_s$ is the Hilbert direct sum of symmetric powers. Mode $j$
corresponds to the normalized Heisenberg creation operator
$a_{-j}^{(p)}/\sqrt j$. Splitting the one-particle space into modes $j=1$
and $j\geq2$ gives
\[
 \mathscr F_F^{\mathrm{cur}}
  \simeq L_F\widehat\otimes\mathscr F_F^{\geq2},
 \qquad
 L_F=\widehat\bigotimes_{p\in F}\mathscr F_p.
\]
The orthonormal number-state bases prove this unitary directly, and it
preserves the grading assigning degree $j$ to a mode-$j$ creation.

Let $\mathbb N_F$ be the integers supported on $F$. The phases and shifts
$\mu_n$, $n\in\mathbb N_F$, act on $L_F\simeq\ell^2(\mathbb N_F)$ and extend
by the identity on $\mathscr F_F^{\geq2}$. Every closed subspace
$L_F\widehat\otimes K_0$, with $K_0\subset\mathscr F_F^{\geq2}$, is
invariant under these operators and their adjoints. If $F$ is nonempty, the
higher-mode vacuum is properly contained in larger such stable sectors.

This restricted representation need not be faithful on all phases. If
$q\notin F$, the nonzero projection
\[
 e_q=\frac1q\sum_{j=0}^{q-1}e(j/q)
\]
acts as zero on $L_F$, since no integer in $\mathbb N_F$ is divisible by
$q$. Also the higher-mode factor by itself is not the full vacuum Heisenberg
vertex algebra: it omits every mode-one creation.

\subsection{The Gibbs trace and its domain}

Define the occupation Hamiltonian by
\[
 H\delta_m=(\log m)\delta_m,\qquad
 \operatorname{Dom}H
  =\left\{\xi:\sum_m(\log m)^2|\xi_m|^2<\infty\right\}.
\]
This is a self-adjoint diagonal operator with finite-support core. For a
real $\beta>1$,
\begin{equation}
\label{v4-arith:traces:gibbs-zeta}
 \operatorname{Tr}(e^{-\beta H})
   =\sum_{m\geq1}m^{-\beta}
   =\prod_p(1-p^{-\beta})^{-1}
   =\zeta(\beta).
\end{equation}
The series is positive and convergent. Unique factorization proves the Euler
product by increasing finite prime sets. At $\beta=1$ the trace is the
divergent harmonic series.

With the KMS convention $\varphi(ab)=\varphi(b\sigma_{i\beta}(a))$ for
analytic elements, the Gibbs state is
\[
 \varphi_\beta(a)
   =\zeta(\beta)^{-1}\operatorname{Tr}(a e^{-\beta H}).
\]
Indeed $e^{-\beta H}a=\sigma_{i\beta}(a)e^{-\beta H}$ on analytic normal
terms, and trace cyclicity proves the identity. The full Bost--Connes
classification gives a unique $\beta$-KMS state for $0<\beta\leq1$. For
$\beta>1$ its extremal states are indexed by embeddings of the cyclotomic
field into $\mathbb C$. These are Bost--Connes, Theorems 5 and 25, with the
critical case in Section 7. At the critical temperature, existence of a KMS
state does not give a trace-class density in the occupation representation.

The spherical $GL_2(\mathbb Q_p)$ measure and the Jordan-quiver groupoid
mass have their own constructions in Chapters~\ref{v4-arith:hecke:chapter}
and~\ref{v4-arith:hall:chapter}. Relating those objects to the Bost--Connes
system requires an operator or categorical comparison preserving their
states and normalizations. Their prime indexing supplies no such map.

\section{Gauge classes in Hochschild cohomology}

We now use cochains with regular coefficients on $A_{\mathrm{alg}}$. These
compute a cohomology algebra, whereas the earlier cyclic chains compute a
trace complex. All cochains in this first calculation are algebraic $\mathbb
C$-multilinear maps.

Let
\[
 W=\operatorname{Hom}_{\mathbb Z}
        (\mathbb Q_{>0}^{\times},(\mathbb C,+))
   =\prod_p\mathbb C.
\]
A family $\lambda=(\lambda_p)_p$ evaluates on a positive rational by
$\lambda(g)=\sum_p v_p(g)\lambda_p$. Every such sum has finite support, even
if the family does not.

\begin{proposition}[Algebraic gauge derivations]
\label{v4-arith:traces:gauge-derivations}
The formula
\[
 D_\lambda(\mu_re(\gamma)\mu_s^*)
  =\lambda(r/s)\mu_re(\gamma)\mu_s^*
\]
defines a derivation. The map $\lambda\mapsto[D_\lambda]$ injects $W$ into
$HH^1_{\mathrm{alg}}(A_{\mathrm{alg}},A_{\mathrm{alg}})$. Every nonzero
$D_\lambda$ is unbounded in the inherited $C^*$-norm and has no bounded
spatial implementer on $\mathcal L$. Moreover
$HH^0_{\mathrm{alg}}(A_{\mathrm{alg}},A_{\mathrm{alg}}) =\mathbb C$.
\end{proposition}
\begin{proof}
Gauge averaging against characters decomposes the normal algebra as
\[
 A_{\mathrm{alg}}
  =\bigoplus_{g\in\mathbb Q_{>0}^{\times}}(A_{\mathrm{alg}})_g,
 \qquad
 (A_{\mathrm{alg}})_g(A_{\mathrm{alg}})_h
       \subseteq(A_{\mathrm{alg}})_{gh}.
\]
This uses homogeneous components, and does not assert independence of normal
terms of the same degree. Multiplying the component of degree $g$ by
$\lambda(g)$ is consequently well defined. Since
$\lambda(gh)=\lambda(g)+\lambda(h)$, it is a derivation.

One-cocycles are derivations and one-coboundaries are inner derivations. If
$D_\lambda=[b,-]$ for a bounded operator $b$, the norm-one isometries
$\mu_p^k$ give
\[
 k|\lambda_p|=\|D_\lambda(\mu_p^k)\|
      \leq2\|b\|\qquad(k\geq1).
\]
Thus $\lambda_p=0$ for every prime. This also proves injectivity on classes
by taking differences, and proves unboundedness for every nonzero family.

An operator commuting with all phases is diagonal: any two distinct integer
indices are separated by some character $\exp(2\pi i m\gamma)$. Write it as
$T\delta_m=t_m\delta_m$. Commutation with every shift requires $t_{nm}=t_m$,
hence $t_n=t_1$. The represented commutant is scalar. Faithfulness gives the
asserted center and degree-zero cohomology.
\end{proof}

The derivation preserves adjoint exactly when every $\lambda_p$ is purely
imaginary. This follows by comparing its values on $\mu_p$ and $\mu_p^*$.
The product $\prod_p\mathbb C$ parametrizes algebraic derivations. It is not
a convergent sum of bounded operators.

\subsection{Finite cycles detecting every exterior degree}

There is a character $\varepsilon:A_{\mathrm{BC}}\to\mathbb C$ defined by
\[
 \varepsilon(\mu_n)=\varepsilon(\mu_n^*)
  =\varepsilon(e(\gamma))=1.
\]
The presentation verifies every relation, including the transfer average of
$n$ ones. This is a norm-one unital star homomorphism. It is different from
the vacuum vector state, which takes $\mu_p$ to zero.

Let $\mathbb C_\varepsilon$ be the bimodule with left and right actions
through $\varepsilon$. Composition of a regular-coefficient cochain with
$\varepsilon$ is a cochain map. For clarity, the two differentials are
\begin{align*}
 (bF)(a_1,\ldots,a_{n+1})={}&
    a_1F(a_2,\ldots,a_{n+1})\\
 &+\sum_{i=1}^n(-1)^i
     F(a_1,\ldots,a_ia_{i+1},\ldots,a_{n+1})\\
 &+(-1)^{n+1}F(a_1,\ldots,a_n)a_{n+1},
\end{align*}
and
\begin{align*}
 b\bigl(c[a_1|\cdots|a_n]\bigr)={}&
    c\varepsilon(a_1)[a_2|\cdots|a_n]\\
 &+\sum_{i=1}^{n-1}(-1)^i
       c[a_1|\cdots|a_ia_{i+1}|\cdots|a_n]\\
 &+(-1)^n\varepsilon(a_n)c[a_1|\cdots|a_{n-1}].
\end{align*}
Evaluation makes these formulas dual after changing cochain coefficients by
$\varepsilon$.

For distinct primes $P=(p_1,\ldots,p_n)$, write $s_j=\mu_{p_j}$ and set
\[
 c_P=\sum_{\pi\in S_n}\operatorname{sgn}(\pi)
                 1[s_{\pi(1)}|\cdots|s_{\pi(n)}].
\]

This is a cycle with coefficients in $\mathbb C_\varepsilon$. For an
interior face, exchange the adjacent multiplied factors: they commute and
their permutation signs are opposite. For the outer faces, pair a word with
its cyclic rotation. Rotation has sign $(-1)^{n-1}$ and the last face has
sign $(-1)^n$, so those terms cancel as well. For $n=1$ the boundary is
directly zero.

Therefore
\[
 L_P(F)=\sum_{\pi\in S_n}\operatorname{sgn}(\pi)
       \varepsilon\bigl(F(s_{\pi(1)},\ldots,s_{\pi(n)})\bigr)
\]
vanishes on every regular-coefficient coboundary.

\begin{theorem}[An exterior algebra of gauge classes]
\label{v4-arith:traces:exterior-gauge}
Cup product defines an injective graded-algebra map
\[
 \Lambda_{\mathbb C}^*W
   \longhookrightarrow
 HH^*_{\mathrm{alg}}(A_{\mathrm{alg}},A_{\mathrm{alg}}),
 \qquad \lambda\longmapsto[D_\lambda].
\]
The exterior algebra consists of finite sums of finite wedges. It has no
completion in this statement.
\end{theorem}
\begin{proof}
The cup of $n$ derivations is the cochain
\[
 F_{\lambda_1,\ldots,\lambda_n}(a_1,\ldots,a_n)
   =D_{\lambda_1}(a_1)\cdots D_{\lambda_n}(a_n).
\]
The cochain differential obeys the graded product rule, as follows by
cancellation of its interior multiplication terms. Hence this cup is a
cocycle. For two derivations $D,E$, expansion of the Leibniz rules gives
\[
 b(DE)(a,b)=-D(a)E(b)-E(a)D(b).
\]
Their classes anticommute, and their squares vanish over $\mathbb C$. This
constructs the exterior-algebra map.

The finite cycle evaluation is exactly
\[
 L_P(F_{\lambda_1,\ldots,\lambda_n})
   =\det\bigl(\lambda_i(p_j)\bigr)_{1\leq i,j\leq n}.
\]

Every exterior element belongs to $\Lambda^*U$ for some finite-dimensional
subspace $U\subset W$. Prime-coordinate functionals separate points of $U$,
so they span $U^*$. Otherwise their span would have a nonzero annihilator
whose every prime coordinate vanished. A finite selection already spans
$U^*$, and its exterior products detect every nonzero element of
$\Lambda^nU$.

The corresponding $L_P$ proves injectivity in degree $n$. The character
detects degree zero.
\end{proof}

For two distinct primes, the cycle is $1[\mu_p|\mu_q]-1[\mu_q|\mu_p]$. Its
boundary is zero term by term after commuting the shifts, and
$L_{(p,q)}(D_p\smile D_q)=1$, where $D_p$ is the coordinate derivation. For
three primes the six-term alternating cycle gives $L_{(p,q,r)}(D_p\smile
D_q\smile D_r)=1$. Thus these lower bounds are nonzero in every degree. They
are not a classification of all Hochschild classes.

\section{Smooth algebras and the time generator}

The topological domain changes the gauge calculation. For a finite prime set
$F$, let $B_F$ consist of the elements whose orbit under the finite torus
$\mathbb T^F\subset K$ is norm-smooth. Angular derivatives are denoted
$\delta_p$, so $\delta_p(\mu_q)=i\delta_{pq}\mu_q$ and $D_p=\delta_p/i$. For
a nonnegative multiindex $\alpha$ supported in $F$, put
\[
 q_\alpha(a)=\|\delta^\alpha a\|,\qquad
 q_0(a)=\|a\|.
\]
Define
\[
 B_{\mathrm g}=\bigcap_{F\ \mathrm{finite}}B_F
\]
with all these finite-support seminorms. This uses smoothness on ordinary
finite-dimensional tori.

\begin{proposition}[Smooth gauge cohomology]
\label{v4-arith:traces:gauge-smooth}
The spaces $B_F$ and $B_{\mathrm g}$ are Fr\'echet star algebras, norm-dense
in $A_{\mathrm{BC}}$. For jointly continuous multilinear cochains with
regular coefficients, there are injections
\[
 \Lambda^*\left(\bigoplus_{p\in F}\mathbb C D_p\right)
       \longhookrightarrow HH^*_{\mathrm{cont}}(B_F,B_F),
 \qquad
 \Lambda^*\left(\bigoplus_p\mathbb C D_p\right)
       \longhookrightarrow HH^*_{\mathrm{cont}}(B_{\mathrm g},B_{\mathrm g}).
\]
\end{proposition}
\begin{proof}
For $B_F$, a Cauchy sequence in every seminorm gives orbit maps that are
Cauchy in every uniform derivative norm. Smooth Banach-valued functions on a
compact torus are complete. The fundamental theorem of calculus identifies
the uniform limit of first derivatives with the derivative of the uniform
limit. Iteration handles every order. The limit orbit is $\alpha_z(a)$ for
the norm limit $a$. Thus $B_F$ is complete.

There are countably many finite-support multiindices. A Cauchy sequence for
$B_{\mathrm g}$ converges in every $B_F$, and all limits agree with its norm
limit. This proves completeness and metrizability. The iterated product rule
\[
 \delta^\alpha(ab)=
   \sum_{\beta\leq\alpha}\binom{\alpha}{\beta}
       (\delta^\beta a)(\delta^{\alpha-\beta}b)
\]
proves continuous multiplication. Real angular derivatives commute with
adjoint. Normal sums lie in every smooth algebra and are norm-dense.

Each permitted $D_p$ is continuous, with its seminorm bounded by the
seminorm having one additional $p$ derivative. For $B_{\mathrm g}$ this
follows by differentiating on $F\cup\{p\}$ before restricting to $F$. Finite
cup cochains are consequently continuous. The character is continuous
because $|\varepsilon(a)|\leq q_0(a)$. The finite cycles and determinant
evaluations of Theorem~\ref{v4-arith:traces:exterior-gauge} therefore detect
these continuous cohomology classes as well.
\end{proof}

\begin{proposition}[The continuity criterion]
\label{v4-arith:traces:gauge-continuity}
The algebraic $D_\lambda$ extends to a continuous linear map $B_{\mathrm
g}\to A_{\mathrm{BC}}$ exactly when $\lambda$ has finite support. It then
extends as a continuous derivation $B_{\mathrm g}\to B_{\mathrm g}$. On
$B_F$ the criterion is support contained in $F$.
\end{proposition}
\begin{proof}
Continuity would give
\[
 \|D_\lambda(a)\|
   \leq C\max_{\alpha\in E}q_\alpha(a)
\]
for some finite set $E$ of multiindices, including zero. Let $F_0$ contain
their supports. For $q\notin F_0$, every positive derivative on the right
vanishes on $\mu_q^k$, while its norm is one. Hence $k|\lambda_q|\leq C$ for
all $k$, forcing $\lambda_q=0$. On $B_F$ the same argument excludes every
$q\notin F$. Conversely a finite linear combination of the permitted
continuous derivatives has the required extension.
\end{proof}

This criterion concerns this particular family of derivations. It does not
classify every continuous derivation of these algebras.

It also separates gauge smoothness from physical-time smoothness. Enumerate
distinct primes $p_j$ and form
\[
 a=\sum_{j\geq1}2^{-j}\mu_{p_j}^{\,2^j}.
\]
The series converges in norm. Only finitely many terms move under any fixed
finite prime torus, so $a\in B_{\mathrm g}$.

If its time derivative existed in norm, applying it to $\delta_1$ would give
coordinate $i\log p_j$ at each distinct basis vector $\delta_{p_j^{2^j}}$.
Those coordinates are not square summable. Thus $a$ is not even
norm-differentiable under $\sigma$.

\subsection{Closed differentiation in physical time}

Let $\delta$ be the norm generator of $\sigma$:
\[
 \operatorname{Dom}\delta
  =\left\{a:\lim_{t\to0}\frac{\sigma_t(a)-a}{t}
                         \text{ exists in norm}\right\}.
\]
Every normal term is entire analytic, and
\[
 \delta(\mu_re(\gamma)\mu_s^*)
   =i\log(r/s)\mu_re(\gamma)\mu_s^*.
\]
The domain is dense and proper. For a direct one-prime example of
properness, use $\sum_{j\geq1}2^{-j}\mu_p^{2^j}$. The derivative on
$\delta_1$ would have the same nonzero coordinate $i\log p$ at infinitely
many states.

The generator is closed. For $a$ in its domain, differentiating the orbit
gives the Bochner integral identity
\[
 \sigma_t(a)-a=\int_0^t\sigma_s(\delta a)\,ds.
\]
If $a_n\to a$ and $\delta a_n\to b$ in norm, isometry permits passage to the
limit. Division by $t$ and norm continuity then give $\delta a=b$.

The norm-smooth orbit algebra is
\[
 A_\sigma^\infty=\bigcap_{j\geq0}\operatorname{Dom}\delta^j,
 \qquad
 p_k(a)=\sum_{j=0}^k\|\delta^j a\|.
\]
Repeated differentiation identifies this intersection with smoothness of the
orbit. Closedness of $\delta$, applied successively to a Cauchy sequence of
derivatives, proves completeness. The iterated Leibniz rule proves
continuous multiplication, and adjoint is continuous. Thus it is a dense
Fr\'echet star algebra. The estimate $p_k(\delta a)\leq p_{k+1}(a)$ makes
$\delta$ continuous on it.

The class of $\delta$ in
$HH^1_{\mathrm{cont}}(A_\sigma^\infty,A_\sigma^\infty)$ is nonzero. An inner
implementer would be a bounded element of $A_{\mathrm{BC}}$, contrary to the
prime-power estimate in Proposition~\ref{v4-arith:traces:gauge-derivations}.
Here the algebraic family is $\lambda_p=i\log p$. Its infinite support also
explains why it does not extend continuously from the gauge-smooth topology.

In the occupation representation it has an unbounded spatial implementer.
Normal terms preserve $\operatorname{Dom}H$, since their nonzero action
changes the logarithmic weight by the constant $\log(r/s)$. On that domain,
\[
 [H,\mu_re(\gamma)\mu_s^*]
    =\log(r/s)\mu_re(\gamma)\mu_s^*,\qquad
 \delta=i[H,-].
\]
Truncation in the weighted norm proves the finite-support core assertion for
$H$ and justifies passage from basis vectors to its maximal domain. The
scalar commutant computed above and the bicommutant theorem give
$A_{\mathrm{BC}}''=B(\mathcal L)$ in this representation. A GNS
representation from a specified KMS state is another carrier and need not
have this weak closure.

\subsection{Logarithmic dilation and an actual circle quotient}

Logarithmic coordinates give a useful comparison of continuous actions. The
map
\[
 Uf(t)=f(e^t):L^2(\mathbb R_{>0},dr/r)
           \xrightarrow{\sim}L^2(\mathbb R,dt)
\]
is unitary. For $D_af(r)=f(e^{-a}r)$, substitution gives
$UD_aU^{-1}F(t)=F(t-a)$.

Fix $L>0$. The quotient group $\mathbb R_{>0}/e^{L\mathbb Z}$ instead
becomes $\mathbb R/L\mathbb Z$. With measure $dt$ on $[0,L)$, its
orthonormal Fourier basis is
\[
 L^{-1/2}e^{2\pi int/L},\qquad n\in\mathbb Z.
\]
Inverse translation has eigenvalue $e^{-2\pi ina/L}$ on this basis. The
logarithm gives the quotient isomorphism, and Fourier orthogonality proves
the normalization.

This circle Hilbert space is not a nonzero periodic subspace of $L^2(\mathbb
R,dt)$. The squared norm of a periodic function is the sum of its identical
integrals over disjoint periods. It is finite only when the function is zero
almost everywhere.

A period $L$ for Bost--Connes time would require $L\log n\in2\pi\mathbb Z$
for every surviving shift. It cannot hold for both $n=2$ and $n=3$.
Otherwise $\log2/\log3$ would be rational, implying $2^a=3^b$ for positive
integers $a,b$, contrary to unique factorization. A circle KMS comparison
therefore needs its actual quotient or representation, periodic dynamics,
state and intertwiner. The logarithmic coordinate alone does not construct
these data.

\begin{question}[The arithmetic Hochschild comparison]
\label{v4-arith:traces:bc-comparison}
Specify an arithmetic category or algebra, its coefficient ring, trace
bimodule, grading, topology and completion. Construct its comparison with
the Bost--Connes normal algebra or one of the stated smooth algebras.
Determine which gauge classes survive and whether the arithmetic theory has
an integral degree-one carrier generated by logarithmic prime data. An
integral lattice, a degree bound, or a map to a Witt cotangent complex
requires its own construction.
\end{question}

For a Hilbert realization, fix a pre-Hilbert space $L_{\mathrm{alg}}$ and an
arithmetic trace complex $(C_\bullet,b)$. Give a cycle map
$J:L_{\mathrm{alg}}\to Z_d(C_\bullet)$ and a positive semidefinite form $q$
on $H_d(C_\bullet)$. If
\[
 q([Jv],[Jw])=\langle v,w\rangle,
\]
then $v\mapsto[Jv]$ induces an isometry from the completion of
$L_{\mathrm{alg}}$ into the completion of $H_d(C_\bullet)/\ker q$. Density
is needed for surjectivity onto a proposed sector. The action on homology
and its intertwining with $J$ are additional data. The occupation
representation by itself does not supply them.

\section{Explicit positive local pairings}

A bilinear evaluation becomes a Hermitian form only after specifying an
antilinear map in one slot. We use inner products linear in the first
variable throughout this section. The following models calculate the map,
its normalization, and its domain.

\subsection{The Segal form and an archimedean line}

Let $\mathscr F_{\mathrm{alg}}$ be the Heisenberg vacuum module generated by
$\alpha_{-n}$, $n\geq1$, with
\[
 [\alpha_m,\alpha_n]=m\delta_{m+n,0},\qquad
 \alpha_n\Omega=0\quad(n>0).
\]
Normalize $H_\infty(\Omega,\Omega)=1$ and impose
$\alpha_{-n}^\dagger=\alpha_n$. For a finite occupation array $\mathbf
r=(r_1,r_2,\ldots)$, put $v_{\mathbf r}=\prod_n\alpha_{-n}^{r_n}\Omega$.
Commuting each annihilation operator to the right gives
\begin{equation}
\label{v4-arith:traces:segal-gram}
 H_\infty(v_{\mathbf r},v_{\mathbf s})
  =\delta_{\mathbf r,\mathbf s}\prod_n n^{r_n}r_n!.
\end{equation}

The factorial counts repeated contractions of the same mode. The factor $n$
comes from its commutator. Equivalently, for two monomials of lengths $k,l$,
\[
 H_\infty(\alpha_{-n_1}\cdots\alpha_{-n_k}\Omega,
          \alpha_{-m_1}\cdots\alpha_{-m_l}\Omega)
  =\delta_{kl}\sum_{\pi\in S_k}
                   \prod_{i=1}^k n_i\delta_{n_i,m_{\pi(i)}}.
\]
This defines a positive form and its Segal Hilbert completion.

Let $J_\infty$ conjugate coefficients in the monomial basis. The Gram matrix
is real, so $J_\infty$ is an isometric antilinear involution. The bilinear
form
\[
 B_\infty(v,w)=H_\infty(v,J_\infty w)
\]
is bounded by $\|v\|\|w\|$, extends to the completion, and satisfies
$B_\infty(v,J_\infty w)=H_\infty(v,w)$. The complex-linear identity cannot
replace $J_\infty$: on the vacuum, $B_\infty(\Omega,i\Omega)=i$ while
$H_\infty(\Omega,i\Omega)=-i$.

There is also an explicit archimedean comparison of lines. On $\mathbb
R^\times$ use measure $da/|a|$ and set
\[
 W_0(a)=|a|^{1/2}K_0(2\pi|a|),\qquad
 K_0(x)=\int_0^\infty e^{-x\cosh u}\,du\quad(x>0).
\]
This is the parameter-zero spherical Whittaker function in the chosen
normalization. Its squared norm is $\pi/4$. To see the constant, Tonelli
gives
\[
 \int_0^\infty K_0(x)^2\,dx
  =\int_0^\infty\int_0^\infty
                \frac{du\,dv}{\cosh u+\cosh v}.
\]

Put $a=(u+v)/2$ and $b=(u-v)/2$. The Jacobian is two, the region is
$a\geq0$, $|b|\leq a$, and the denominator is $2\cosh a\cosh b$. With
$F(a)=\int_0^a\operatorname{sech}b\,db$, the integral is
\[
 2\int_0^\infty F'(a)F(a)\,da
   =F(\infty)^2=\frac{\pi^2}{4}.
\]

Scaling $x=2\pi|a|$ and including both signs gives $\|W_0\|^2=\pi/4$.

Hence
\[
 \mathbb C\Omega\longrightarrow\mathbb C W_0,\qquad
 z\Omega\longmapsto\frac{2z}{\sqrt\pi}W_0
\]
is an isometry. Extending it to a full Whittaker or arithmetic
representation requires an intertwiner for the specified actions. A line
isometry does not determine one.

\subsection{Weighted occupations and two isometries}

Fix a prime $p$, and define
\[
 H_{\mathrm{occ},p}=\ell^2(\mathbb N_0,p^{-j}),\qquad
 H_{\mathrm{occ}}(v,w)=\sum_{j\geq0}p^{-j}v_j\overline{w_j}.
\]
Its coordinate vector $\delta_j$ has squared norm $p^{-j}$. The raw shift
$R\delta_j=\delta_{j+1}$ has norm $p^{-1/2}$, and $\sqrt p\,R$ is an
isometry. Coordinate conjugation $J_{\mathrm{occ}}$ gives the continuous
bilinear form
\[
 B_{\mathrm{occ}}(v,w)=\sum_{j\geq0}p^{-j}v_jw_j,\qquad
 B_{\mathrm{occ}}(v,J_{\mathrm{occ}}w)=H_{\mathrm{occ}}(v,w).
\]
Cauchy--Schwarz proves absolute convergence and continuity.

Let $E_{ij}$ be the matrix units on $\ell^2(\mathbb N_0)$. Choose a
multiplicative Haar measure on $\mathbb Q_p^\times$ and put
$m_p=\operatorname{vol}(\mathbb Z_p^\times)>0$. There are isometries
\begin{equation}
\label{v4-arith:traces:local-isometries}
 D\delta_j=p^{-j/2}E_{jj},\qquad
 T\delta_j=p^{-j/2}m_p^{-1/2}\mathbf1_{p^j\mathbb Z_p^\times}.
\end{equation}
Their ranges are the diagonal Hilbert--Schmidt subspace and the closed
radial-function subspace supported on nonnegative valuations, respectively.

Indeed matrix units are Hilbert--Schmidt orthonormal. The disjoint shells
$p^j\mathbb Z_p^\times$ each have multiplicative measure $m_p$. Both
displayed families therefore have Gram matrix $p^{-j}\delta_{jk}$. Finite
spans are dense in the stated ranges, proving the isometric extensions.

Their basis vectors are fixed by adjunction or function conjugation. Thus
\[
 H_{\mathrm{occ}}(v,w)
  =\operatorname{Tr}(Dv(Dw)^*)
  =\int_{\mathbb Q_p^\times}Tv(a)\overline{Tw(a)}\,d^\times a.
\]
Also $B_{\mathrm{occ}}(v,w)=\operatorname{Tr}(DvDw)$.

The diagonal Hilbert--Schmidt range is not invariant under left
multiplication by the shift. For example, $VE_{jj}=E_{j+1,j}$. Consequently
this isometry does not identify the whole prime GNS representation. Nor can
its dual basis be realized as $\mu_{p^j}^*\delta_1$ in the occupation
representation: those vectors vanish for $j>0$.

\subsection{The spherical norm depends on the representation}

Fix a normalized unramified tempered representation $\pi_p$ of $GL_2(\mathbb
Q_p)$, with Satake parameters $|\alpha|=|\beta|=1$ and an additive character
of conductor $\mathbb Z_p$. Its spherical Whittaker function normalized by
$W(1)=1$ has shell values
\[
 W(\operatorname{diag}(p^j,1))
   =p^{-j/2}h_j(\alpha,\beta),\qquad
 h_j(\alpha,\beta)=\sum_{r=0}^j\alpha^{j-r}\beta^r
 \quad(j\geq0),
\]
and vanishes at negative valuations. We use this standard spherical
Whittaker formula with the stated normalization.

\begin{proposition}[The full spherical norm]
\label{v4-arith:traces:spherical-norm}
For $\operatorname{Re}s>0$,
\[
 Z_p(s)=\int_{\mathbb Q_p^\times}|W(a)|^2|a|^{s-1}\,d^\times a
   =m_p\frac{L_p(s,\pi_p\times\widetilde\pi_p)}{\zeta_p(2s)}.
\]
In particular,
\[
 C_p=\|W\|^2
   =m_p(1+p^{-1})L_p(1,\operatorname{Ad}\pi_p)>0.
\]
\end{proposition}
\begin{proof}
Shell integration gives $Z_p(s)=m_p\sum_{j\geq0}|h_j|^2p^{-js}$. The bound
$|h_j|\leq j+1$ proves absolute convergence on the indicated half-plane. Put
$r=\alpha/\beta$. For $|x|<1$, geometric summation gives
\[
 \sum_{j\geq0}|h_j|^2x^j
   =\frac{1+x}{(1-x)(1-rx)(1-r^{-1}x)}.
\]
For $r\ne1$, expand $|h_j|^2=(2-r^{j+1}-r^{-j-1})/|1-r|^2$ and sum its three
geometric series. For $r=1$, use $\sum_{j\geq0}(j+1)^2x^j=(1+x)/(1-x)^3$.

The Rankin--Selberg denominator is $(1-x)^2(1-rx)(1-r^{-1}x)$.
Multiplication by $1-x^2$ proves the formula with $x=p^{-s}$. Since
\[
 L_p(s,\pi_p\times\widetilde\pi_p)
  =\zeta_p(s)L_p(s,\operatorname{Ad}\pi_p),
\]
evaluation at one gives $C_p$. It is a positive norm, and its denominator
factors are nonzero because $p^{-1}<1$.
\end{proof}

For $d^\times a=da/|a|$ with additive unit volume, $m_p=1-p^{-1}$ and
$C_p=(1-p^{-2})L_p(1,\operatorname{Ad}\pi_p)$. In contrast, the compact
subgroup volume for the invariant matrix measure $|\det g|^{-2}\prod
dg_{ij}$ is
\[
 \operatorname{vol}(GL_2(\mathbb Z_p))
    =\frac{(p^2-1)(p^2-p)}{p^4}
    =(1-p^{-1})(1-p^{-2}).
\]
Reduction modulo $p$ proves this formula by counting the two independent
columns. It is not the Whittaker norm, which depends on $\alpha/\beta$.
Inserting an additional $|a|^{-1}$ into the norm integral changes the
question to $Z_p(0)$. For $\alpha=\beta=1$, the resulting shell series is
$m_p\sum_{j\geq0}(j+1)^2$ and diverges.

The radial isometry of \eqref{v4-arith:traces:local-isometries} identifies
\[
 W=T\left(\sqrt{m_p}\sum_{j\geq0}h_j(\alpha,\beta)\delta_j\right).
\]
This vector converges in $H_{\mathrm{occ},p}$ and has squared norm $C_p$. It
is generally different from $T\delta_0$. Unit normalization uses
$W/\sqrt{C_p}$. Dividing a pairing by the adjoint factor alone gives
spherical squared norm $m_p(1+p^{-1})$, rather than one.

\section{The prime Gibbs state and its modular conjugation}

The full prime cyclic space can also be constructed explicitly. Let
$K_0=\ell^2(\mathbb N_0)$, $Ve_j=e_{j+1}$, and $H_pe_j=j\log p\,e_j$, with
maximal domain $\{\xi:\sum j^2|\xi_j|^2<\infty\}$. For $\beta>0$, put
\[
 q=p^{-\beta},\qquad
 r_j=(1-q)q^j,\qquad
 \rho e_j=r_je_j.
\]
The positive injective operator $\rho$ has trace one. Its state on the
Toeplitz algebra $\mathcal T=C^*(V)$ is
\[
 \varphi_{p,\beta}(a)=\operatorname{Tr}(\rho a),\qquad
 \varphi_{p,\beta}(V^mV^{*n})=\delta_{mn}q^m.
\]
The physical time is $\alpha_t(a)=e^{itH_p}ae^{-itH_p}$, with
$\alpha_t(V)=p^{it}V$.

This is the unique $\beta$-KMS state on $(\mathcal T,\alpha)$. On analytic
monomials, $\rho a=\alpha_{i\beta}(a)\rho$, so the KMS identity follows by
trace cyclicity. Conversely that identity with second entry one kills each
monomial of nonzero degree. With first entry $V^m$ and second entry
$V^{*m}$, it gives the value $q^m$. Density of the analytic monomials proves
uniqueness.

\begin{proposition}[The Hilbert--Schmidt GNS model]
\label{v4-arith:traces:prime-gns}
The GNS space and representation are
\[
 \mathcal H_{p,\beta}=\operatorname{HS}(K_0),\qquad
 \pi(a)X=aX,\qquad
 \Omega=\rho^{1/2},
\]
with inner product $\langle X,Y\rangle=\operatorname{Tr}(XY^*)$. The GNS map
is $[a]\mapsto a\rho^{1/2}$. The vector $\Omega$ is cyclic and separating
for left $B(K_0)=\pi(\mathcal T)''$.
\end{proposition}
\begin{proof}
The map is isometric because its inner product is $\operatorname{Tr}(\rho
b^*a)=\varphi_{p,\beta}(b^*a)$. The matrix units $E_{ij}=V^i(1-VV^*)V^{*j}$
belong to $\mathcal T$, and $E_{ij}\rho^{1/2}=\sqrt{r_j}E_{ij}$. They span a
dense subspace of $\operatorname{HS}(K_0)$. Finite matrix corners strongly
approximate bounded operators, proving the weak-closure assertion. If
$a\Omega=0$, every column of $a$ vanishes because $r_j>0$ for all $j$. Thus
$\Omega$ is separating.
\end{proof}

Define the antilinear operator $S_0(a\Omega)=a^*\Omega$ on $\mathcal
T\Omega$.

\begin{proposition}[Closed modular operators and their domains]
\label{v4-arith:traces:prime-modular}
Set $JX=X^*$. For $X=\sum x_{ij}E_{ij}$, the positive modular operator
satisfies
\[
 \Delta E_{ij}=q^{i-j}E_{ij},\qquad
 \operatorname{Dom}\Delta^a
  =\left\{X:\sum_{i,j}q^{2a(i-j)}|x_{ij}|^2<\infty\right\}
       \quad(a\in\mathbb R).
\]
One has $\overline{S_0}=J\Delta^{1/2}$ and
\[
 \Delta^{it}X=e^{-i\beta tH_p}Xe^{i\beta tH_p}.
\]
The modular action on left multiplication is physical time $\alpha_{-\beta
t}$. Both $\Delta$ and $\Delta^{-1}$ are unbounded.
\end{proposition}
\begin{proof}
On matrix units, $S_0E_{ij}=q^{(i-j)/2}E_{ji}$ antilinearly. Truncation of a
convergent weighted double sum makes finite matrix sums graph-dense in the
displayed domain for $a=1/2$. For bounded $a$, the vector $a\Omega$ belongs
to that domain and the formula gives $a^*\Omega$. It follows that the
closure is precisely $J\Delta^{1/2}$. Diagonal spectral calculus proves all
powers and the unitary group. Its eigenvalues range through $q^{\mathbb Z}$,
proving both unboundedness assertions.
\end{proof}

The logarithm has eigenvalue $(i-j)\log q$ on $E_{ij}$ and maximal domain
$\{X:\sum |(i-j)\log q|^2|x_{ij}|^2<\infty\}$. The conjugation acts on
prime-shift vectors by
\[
 J(V^k\Omega)=p^{\beta k/2}V^{*k}\Omega,\qquad
 J(V^{*k}\Omega)=p^{-\beta k/2}V^k\Omega.
\]
Expand $V^k\Omega=\sum_j\sqrt{r_j}E_{j+k,j}$ and use $r_{j+k}=q^kr_j$ to
verify the factors. In particular $\|V^k\Omega\|=1$ and
$\|V^{*k}\Omega\|=q^{k/2}$. The algebra star operation alone is not this
antiunitary on the cyclic space.

For left and right multiplication, $JL_aJ=R_{a^*}$. The continuous bilinear
form $B(X,Z)=\operatorname{Tr}(XZ)$ satisfies
\[
 B(X,JY)=\operatorname{Tr}(XY^*)=\langle X,Y\rangle.
\]
The product of two Hilbert--Schmidt operators is trace class and
$|\operatorname{Tr}(XZ)|\leq\|X\|_2\|Z\|_2$, which proves the asserted
continuity.

\begin{proposition}[Comparison with a specified prime cyclic space]
\label{v4-arith:traces:prime-cyclic-comparison}
Let $(\mathcal A,\alpha,\psi_\beta)$ be a dynamical $C^*$-algebra with a
$\beta$-KMS state and an isometry $\mu_p$ satisfying
$\alpha_t(\mu_p)=p^{it}\mu_p$. The map $\iota_p:\mathcal T\to\mathcal A$
sending $V$ to $\mu_p$ induces a unitary
\[
 U_p:\operatorname{HS}(K_0)
   \longrightarrow
   \overline{\pi_{\psi_\beta}(C^*(\mu_p))\Omega_{\psi_\beta}},
 \qquad
 a\rho^{1/2}\longmapsto
 \pi_{\psi_\beta}(\iota_p(a))\Omega_{\psi_\beta}.
\]
It preserves the prime algebra action, cyclic vector, and closed modular
operators of this restricted cyclic representation, including when
$\beta=1$.
\end{proposition}
\begin{proof}
The pulled-back state is a KMS state on $\mathcal T$, so uniqueness
identifies it with $\varphi_{p,\beta}$. The displayed map is therefore
isometric on $\mathcal T\rho^{1/2}$ and has dense image in the specified
cyclic subspace. It extends to a unitary. It intertwines multiplication and
$a\Omega\mapsto a^*\Omega$ on the core, hence their closures and polar
decompositions.
\end{proof}

The target is the prime-generated cyclic subspace, not the whole global GNS
space. Its identification says nothing yet about additional phase operators
or an arithmetic automorphic placement. For finitely many primes, tensoring
the densities and matrix formulas gives their finite Gibbs and GNS products
for every $\beta>0$. The global occupation heat trace still diverges at
$\beta=1$, as \eqref{v4-arith:traces:gibbs-zeta} shows.

\section{Restricted tensor pairings and arithmetic placement}

Choose complex Hilbert spaces $(H_v,e_v)$ indexed by the places, with unit
reference vectors. Suppose each has a conjugation $J_v$, with $J_ve_v=e_v$,
and a bilinear form $B_v(x,J_vy)=H_v(x,y)$. The algebraic restricted tensor
consists of elementary tensors equal to $e_v$ outside a finite set.

\begin{proposition}[Assembly of normalized forms]
\label{v4-arith:traces:pairing-assembly}
There are forms
\[
 H(\otimes'_v x_v,\otimes'_v y_v)
     =\prod_v H_v(x_v,y_v),\qquad
 B(\otimes'_v x_v,\otimes'_v y_v)
     =\prod_v B_v(x_v,y_v),
\]
and a conjugation $J=\otimes'_v J_v$. They extend to the Hilbert restricted
tensor product, where $H(x,y)=B(x,Jy)$. Its reference vector has squared
norm one.
\end{proposition}
\begin{proof}
At each finite set of places, use the Hilbert tensor form. Enlarging the set
tensors with unit reference vectors, so its insertion maps are isometries.
The directed union is a pre-Hilbert space. The local conjugations give
compatible antilinear isometries, whose square is the identity. They extend
to a conjugation on its completion.

Define $B(x,y)=H(x,Jy)$ there. It is bounded and agrees with the displayed
finite products. All omitted reference factors equal one.
\end{proof}

Raw local norms cannot be multiplied without addressing convergence. For the
invariant matrix measure on $GL_n(\mathbb Q_p)$,
\[
 \operatorname{vol}(GL_n(\mathbb Z_p))
     =\prod_{j=1}^n(1-p^{-j}).
\]
Counting independent columns modulo $p$ proves this identity. The product of
its $j=1$ factors tends to zero. The factors with $j\geq2$ multiply to
$\prod_{j=2}^n\zeta(j)^{-1}$.

A Tamagawa measure inserts its convergence factors, global residue
normalization, and rational-character quotient. It is not the raw product of
these volumes. Nor does the volume of an automorphic quotient determine the
Petersson norm of a specified cusp form.

For example, with the Tamagawa measure disintegrated along $|\det|$ using
$dt/t$, the volume of
\[
 GL_2(\mathbb Q)\backslash GL_2(\mathbb A)^1,\qquad
 GL_2(\mathbb A)^1=\ker|\det|_{\mathbb A},
\]
is one. This is the classical Tamagawa number in its rational-character
normalization, as in Lai, Section 1.7 and Theorem 7.1.

In that formula the split diagonal torus has Tamagawa number one.
Restricting its character $(a,b)$ to the scalar torus and the derived
diagonal torus gives $(a+b,a-b)$, a sublattice of index two. The central
characters trivial on their intersection $\mu_2$ are also a sublattice of
index two. The two index factors cancel, giving $\tau(GL_2)=\tau(\mathbb
G_m^2)=1$.

This global volume calculation does not normalize the vector $W$ or an
individual automorphic form.

The local spherical factors are likewise representation dependent. If a
specified global adjoint $L$-function continues holomorphically to one with
$c_\pi=L(1,\operatorname{Ad}\pi)\ne0$, then $c_\pi^{-1}B$ is a well-defined
rescaled pairing. Rescaling both sides of a pairing identity by that same
scalar preserves the identity. It does not prove ordinary convergence of the
local product at one.

For tempered unramified parameters and $\sigma>1$, the adjoint logarithmic
series has the bound
\[
 \sum_p\sum_{j\geq1}\frac{3p^{-j\sigma}}j
   \leq\frac{3}{1-2^{-\sigma}}\sum_p p^{-\sigma}<\infty.
\]
This gives an initial half-plane of absolute convergence. Continuation and
nonvanishing at one are additional analytic inputs for the representation in
question. Finitely many ramified factors require their own normalization and
pole conditions.

\begin{hypothesis}[Arithmetic pairing comparison]
\label{v4-arith:traces:pairing-hypothesis}
Fix an arithmetic state space $V$, a dual coefficient space $V^!$, a
bilinear evaluation $B^{\mathrm{Ar}}:V\otimes V^!\to\mathbb C$, and an
antilinear map $\sigma:V\to V^!$. Specify an injective placement $\iota$ of
a common dense domain into an automorphic Hilbert space with fixed central
character and fixed Haar normalization. At each place require maps on both
pairing carriers, preserving the stated actions. They must carry the
evaluation to a chosen local model and $\sigma$ to its conjugation.

Require the maps to preserve the reference vectors at almost all places and
to identify the placed local Petersson form with the model Hermitian form.
Specify the resulting adelic factorization with a common global scalar.
Finally require continuous extension on the stated completion, or a
specified regularized extension compatible with those maps. Every included
ramified, residual, or continuous sector requires its own local maps and
normalization.
\end{hypothesis}

\begin{proposition}[The conditional Koszul--Petersson identity]
\label{v4-arith:traces:arithmetic-pairing}
Under Hypothesis~\ref{v4-arith:traces:pairing-hypothesis}, on its specified
domain,
\[
 \langle\iota v,\iota w\rangle_{\mathrm{Pet}}
      =B^{\mathrm{Ar}}(v,\sigma w).
\]
If the placement is into an actual $L^2$ space, this transported Hermitian
form is positive definite.
\end{proposition}
\begin{proof}
For elementary restricted tensors, the local model identity
$H_v(x,y)=B_v(x,J_vy)$ gives equality at each place. The reference
normalization and common global scalar make their finite products equal. The
stipulated extension carries the equality to its stated domain. If $v\ne0$,
injectivity gives $\|\iota v\|^2>0$, proving positivity.
\end{proof}

The holomorphic sector supplies concrete local examples. Let
$f=\sum_{n\geq1}a_nq^n$ be a normalized holomorphic cuspidal newform of
weight $k\geq2$, level $N$ and nebentypus $\chi$. At $p\nmid N$, write
\[
 X^2-a_pX+\chi(p)p^{k-1}
   =(X-\alpha_p^{\mathrm{alg}})(X-\beta_p^{\mathrm{alg}}).
\]
Deligne's Theorem 8.2 in \emph{La conjecture de Weil I} gives
$|\alpha_p^{\mathrm{alg}}|=|\beta_p^{\mathrm{alg}}| =p^{(k-1)/2}$. The
unitary parameters $\alpha_p=p^{-(k-1)/2}\alpha_p^{\mathrm{alg}}$ and
$\beta_p=p^{-(k-1)/2}\beta_p^{\mathrm{alg}}$ therefore satisfy the
hypotheses of the spherical norm calculation.

For a cuspidal weight-one eigenform, Deligne--Serre's Theorem 4.1 constructs
the associated continuous complex Galois representation. Its finite-image
Frobenius eigenvalues are roots of unity, giving the same unitary conclusion
at unramified primes. This is their Corollary 4.2. The archimedean
weight-one representation is a limit of discrete series. It is not a
discrete-series representation of weight at least two.

These theorems identify the unramified local parameters. They do not supply
the arithmetic placement, ramified pairing maps or adelic extension in
Hypothesis~\ref{v4-arith:traces:pairing-hypothesis}. Once those data are
given on a selected holomorphic sector, the preceding proposition proves its
pairing identity. The same implication applies to a selected Maass sector
with its own local models and maps.

The positivity implication itself needs no temperedness assumption.
Temperedness enters the separate local spherical model above. Using that
model for a Maass sector requires the relevant local bounds as well as the
arithmetic maps. A local exponent estimate is neither such a map nor an
adelic summability statement. For example $\sum_p p^{-25/64}$ diverges, so
this power cannot be an absolutely summable prime error bound.

The hypothesis preserves the original comparison problem even on a
normalized cuspidal core. An adjunction alone does not identify its counit
target with $\mathbb C$, its dual carrier with Bost--Connes operators, or
its conjugation with an automorphic one. If a placement $U$ intertwines the
model conjugation with function conjugation on a dense domain, continuity
extends that equality to the completion. In a fixed central-character space,
function conjugation naturally lands in the conjugate-character space, which
must be included in the specified target.

Positive forms also do not by themselves make an action unitary. One must
prove the adjoint identities for its operators on invariant domains and
establish any closures used by the representation. The bounded local
isometries above provide the pairings and lines they state. The arithmetic
intertwiners remain to be constructed.

\begin{proposition}[Two dense realizations of one pairing]
\label{v4-arith:traces:dense-pairing}
Let $E$ be a complex vector space with a positive semidefinite Hermitian
form $q$. Suppose linear maps $I:E\to H_1$ and $J:E\to H_2$ have dense
images in Hilbert spaces and satisfy
\[
 q(v,w)=\langle Iv,Iw\rangle_1=\langle Jv,Jw\rangle_2.
\]
Then $Iv\mapsto Jv$ extends uniquely to a unitary $H_1\to H_2$.
\end{proposition}
\begin{proof}
Both kernels are $\{v:q(v,v)=0\}$. The proposed map is consequently well
defined and isometric on the dense image of $I$. Extend by continuity. An
isometry from a complete space has closed range, and this range contains the
dense image of $J$. It is therefore all of $H_2$. Density proves uniqueness.
\end{proof}

For a Petersson--chiral comparison, one must construct $E,B,\sigma,I,J$ with
$q(v,w)=B(v,\sigma w)$ on the stated automorphic and chiral carriers. An
Arnold or Knizhnik--Zamolodchikov connection further requires bundles over
an identified base and a map preserving their connections. A Hilbert-space
isometry does not by itself compare those connections.

\subsection{Meromorphic comparison}

Suppose two scalar pairings of specified meromorphic families of vectors
extend to meromorphic functions on a connected parameter domain. If the
actual comparison identifies them on a nonempty open subset, they agree
everywhere meromorphically. Locally clear their poles and apply the identity
theorem. Their Laurent coefficients, pole orders and residues then agree.

This implication requires both families and the comparison. The local
spherical function $Z_p(s)=m_pL_p(s,\pi_p\times\widetilde\pi_p)/\zeta_p(2s)$
uses a rank-four Rankin--Selberg factor. It is not the rank-two standard
factor $(1-\alpha p^{-s})^{-1}(1-\beta p^{-s})^{-1}$. It does not imply a
universal assertion that every arithmetic pairing has only two Tate boundary
poles. A meromorphic line bundle alone has no pole divisor: that divisor
belongs to a specified meromorphic section in its local trivializations.

\section{Instanton sums, Fock traces, and their comparison}

An equivariant instanton integral and a trace on symmetric powers weight
partitions differently. Their distinction is visible before taking an
arithmetic product.

Let $\lambda=(\lambda_1\geq\lambda_2\geq\cdots)$ be a partition, with
$\lambda_i=0$ beyond its last row. For a box $s=(i,j)$, define
\[
 A_\lambda(s)=\lambda_i-j,\qquad
 L_\lambda(s)=\lambda'_j-i.
\]
For a box outside the diagram these integers can be negative. In particular
the leg of a box in one partition, measured against a different partition,
is not necessarily nonnegative.

For a rank $r$ framing, let
$\boldsymbol\lambda=(\lambda^{(1)},\ldots,\lambda^{(r)})$,
$|\boldsymbol\lambda|=\sum_\alpha|\lambda^{(\alpha)}|$, and choose
equivariant parameters $a_1,\ldots,a_r,\epsilon_1,\epsilon_2$. Write
$\epsilon_+=\epsilon_1+\epsilon_2$ and
\[
 E_{\alpha\beta}(s)
  =a_\alpha-a_\beta-\epsilon_1L_{\lambda^{(\beta)}}(s)
                  +\epsilon_2(A_{\lambda^{(\alpha)}}(s)+1).
\]
The pure vector-multiplet fixed-point sum is
\begin{equation}
\label{v4-arith:traces:instanton}
 Z_{\mathrm{inst}}
   =\sum_{\boldsymbol\lambda}q^{|\boldsymbol\lambda|}
     \prod_{\alpha,\beta=1}^r
     \prod_{s\in\lambda^{(\alpha)}}
       \frac1{E_{\alpha\beta}(s)
                    (\epsilon_+-E_{\alpha\beta}(s))}.
\end{equation}

This is a formal series in the instanton-counting parameter $q$, with
rational equivariant coefficients. Specialization requires its denominators
to be defined. In the pure-gauge scale convention of Nekrasov--Okounkov,
$q=\Lambda^{2r}$. Keeping $q$ explicit avoids confusing that convention with
a change in rank. Matter multiplets contribute additional fixed-point
factors.

The localization construction uses framed sheaves on $\mathbb C^2$, or its
projective compactification with a framing at infinity. This is the
equivariant instanton sum of Nekrasov and Okounkov. For the signed arm
and leg convention, see Alday, Gaiotto and Tachikawa, Appendix B.1.

\begin{proposition}[The fixed-parameter large-equivariant limit]
\label{v4-arith:traces:instanton-limit}
For $r=1$, the one-box coefficient is $1/(\epsilon_1\epsilon_2)$. If
$\epsilon_1\ne0$ and $q$ are fixed while $\epsilon_2\to\infty$, the series
\eqref{v4-arith:traces:instanton} tends coefficientwise in $q$ to $1$.
\end{proposition}
\begin{proof}
For the single box, $A=L=0$ and $E=\epsilon_2$, $\epsilon_+-E=\epsilon_1$.
For a general box in a single partition,
\[
 E=-\epsilon_1L+\epsilon_2(A+1),\qquad
 \epsilon_+-E=\epsilon_1(L+1)-\epsilon_2A.
\]
Here $A,L\geq0$. The first factor has nonzero leading coefficient in
$\epsilon_2$. If $A>0$, the second also grows linearly. If $A=0$, it is the
nonzero constant $\epsilon_1(L+1)$.

Every nonempty partition therefore contributes a coefficient tending to
zero. There are only finitely many partitions of each fixed size, so their
sum tends to zero as well.
\end{proof}

Thus retaining a term $1+q$ requires a different scaling or normalization.
Such a scaling must be specified and tested on all charges. Even an actual
two-term local series would not be the bosonic Euler factor $(1-q)^{-1}$.

\subsection{The hook-weighted specialization}

On the slice $\epsilon_1=-\epsilon_2=\hbar\ne0$, the rank-one denominator of
each box is $-\hbar^2h_\lambda(s)^2$, where
$h_\lambda(s)=A_\lambda(s)+L_\lambda(s)+1$. After setting $x=q/(-\hbar^2)$,
the normalized sum is
\[
 Z_{\mathrm{hook}}(x)
    =\sum_\lambda
        \frac{x^{|\lambda|}}{\prod_{s\in\lambda}h_\lambda(s)^2}.
\]

The hook-length formula for the irreducible $S_n$-representation indexed by
$\lambda$ says $f^\lambda=n!/\prod_s h_\lambda(s)$. Decomposing the regular
representation of $S_n$ gives $\sum_{\lambda\vdash n}(f^\lambda)^2=n!$.
Consequently
\[
 \sum_{\lambda\vdash n}
          \prod_{s\in\lambda}h_\lambda(s)^{-2}
     =\frac1{n!},\qquad
 Z_{\mathrm{hook}}(x)=e^x.
\]

This is the Plancherel specialization in Nekrasov--Okounkov, Section 3.2.
For $n=2$, the two diagrams each contribute $1/4$, so the coefficient is
$1/2$, already different from the coefficient one of $(1-x)^{-1}$.

Substituting $x=p^{-s}$ and multiplying over primes gives $\exp(\sum_p
p^{-s})$ when $\operatorname{Re}s>1$. It retains the first prime-power term
of the logarithm of the Euler product. The higher prime powers are different
coefficients.

\subsection{The all-charge Fock trace}

Let $V$ be finite-dimensional over $\mathbb C$, and let
$T\in\operatorname{End}(V)$. Then
\begin{equation}
\label{v4-arith:traces:fock-determinant}
 \sum_{n\geq0}q^n\operatorname{Tr}(\operatorname{Sym}^nT)
    =\det(1-qT)^{-1}
\end{equation}
as a formal series. Triangularize $T$. On monomials in a triangular basis,
the diagonal entries are products of its eigenvalues. Summing their
nonnegative multiplicities gives the product of geometric series, proving
the identity without a semisimplicity assumption. For the Hilbert direct sum
of symmetric powers, $|q|\|T\|<1$ is a sufficient trace-class condition.

For one oscillator with $T=1$, every charge occurs once, so its trace is
$(1-q)^{-1}$. At $q=p^{-s}$ this is the primon local factor. For a specified
unramified Frobenius operator $F_p$ on a rank-$r$ space, the corresponding
local trace is $\det(1-p^{-s}F_p)^{-1}$. This rank is the dimension of that
coefficient space. It is not inferred from the number of primes in an
indexing set.

For $\operatorname{Re}s>1$, choose the logarithm of $\zeta(s)$ that tends to
zero as $\operatorname{Re}s\to+\infty$. Absolute convergence gives
\begin{equation}
\label{v4-arith:traces:euler-log}
 \log\zeta(s)=\sum_p\sum_{k\geq1}\frac{p^{-ks}}k.
\end{equation}
The bound $\sum_{p,k}p^{-k\sigma}/k \leq(1-2^{-\sigma})^{-1}\sum_p
p^{-\sigma}$ justifies rearrangement. The identical argument applies to the
specified unramified factors of a character in a half-plane of absolute
convergence.

One may define a local scalar $\mathcal W_p(s)=-\log(1-p^{-s})$ by this
formula. Identifying it with a Nekrasov--Shatashvili superpotential requires
an equivariant family and the actual scaled limit defining that
superpotential. The equality \eqref{v4-arith:traces:euler-log} alone does
not construct such a family.

\subsection{The arithmetic instanton comparison}

An arithmetic curve and the framed complex surface used in localization are
different geometric inputs. A proposed arithmetic self-product must specify
its base category, structure sheaf, moduli problem, equivariant action,
fixed loci and integration theory. Listing prime pairs does not construct
those data. Likewise a countable family of local oscillators does not
determine a gauge group or its Coulomb branch.

The real parameter in a metrized arithmetic line can be recorded by its
Arakelov degree. A unitary character supplies separate phase data. This
gives meaningful scalar parameters, but no identification with the Coulomb
differences $a_\alpha-a_\beta$ in \eqref{v4-arith:traces:instanton} until
the comparison maps are supplied. An arithmetic intersection pairing must
first specify its arithmetic surface, metrized bundles, and cycle classes.

\begin{question}[From an arithmetic moduli problem to Euler factors]
\label{v4-arith:traces:arithmetic-instantons}
Construct an arithmetic moduli problem with an equivariant or motivic
integration theory and a comparison to the framed instanton construction.
Specify the coefficient ring, potential or vanishing-cycle data when used,
and every parameter scaling. Determine whether its all-charge trace is the
prescribed Frobenius determinant and prove that comparison on the actual
objects. The fixed-parameter limit and hook-weighted specialization above do
not supply it.
\end{question}

\section{Geometric transfer operators and Liouville scaling}

The geometric $R$-matrix construction contains commuting transfer operators,
but it retains more than their scalar traces. Suppose $R_{F,W}(u)$ is a
compatible family of invertible $R$-matrices satisfying the Yang--Baxter
equation. For a finite-dimensional auxiliary space $F$ and an endomorphism
$g_F$, define
\[
 T_F(g,u)=\operatorname{Tr}_F
                  \bigl((g_F\otimes1)R_{F,W}(u)\bigr).
\]

Assume the family $g$ commutes with the $R$-matrices on each pair of
auxiliary spaces. The Yang--Baxter equation conjugates the two ordered
products by the auxiliary $R$-matrix. Taking the trace over both auxiliary
factors therefore gives
\[
 [T_{F_1}(g,u_1),T_{F_2}(g,u_2)]=0.
\]

For infinite-dimensional auxiliary spaces, the trace needs its specified
grading, completion, or convergence condition.

Maulik--Okounkov construct such operators from stable envelopes of Nakajima
varieties. Their Theorem 1.3.1 identifies the cohomological degree-two part
of the Baxter algebra with modified quantum multiplication by tautological
degree-two classes. The modification includes the theta-characteristic sign
on curve classes. In the divisor formula, the nonclassical part has the form
\[
 -\hbar\sum_{\theta\cdot\alpha>0}
     (\lambda,\alpha)\frac{q^\alpha}{1-q^\alpha}
                   e_\alpha e_{-\alpha},
\]
with dual root-space bases and the scalar term fixed by the unit condition.
This is a geometric theorem on its equivariant cohomology carrier.

For the Jordan quiver, with one vertex and one loop, the variety $\mathcal
M(r,n)$ parametrizes framed rank-$r$ torsion-free sheaves on $\mathbb P^2$
with second Chern class $n$. Its Yangian belongs to the affine
$\widehat{\mathfrak{gl}}_1$ construction. The integer $r$ is the framing
rank and labels the representation. It does not replace the quiver by affine
$\mathfrak{gl}_r$. At $r=1$, the Hilbert-scheme cohomology has its
Heisenberg Fock realization. Stable envelopes relate the rank-$r$ cohomology
to $r$ such bosonic factors on the equivariant domain.

A comparison with arithmetic Hecke operators must specify the geometric
representation, arithmetic representation, parameter map and intertwiner.
Commutativity of a Baxter algebra does not erase the $R$-matrix or its
tensor structure. Identifying a scalar instanton partition sum with a
holonomy trace requires another comparison. These are additional equalities
to construct.

\subsection{The semiclassical Liouville equation}

For Liouville theory with a reference metric $\widetilde g$, use
\[
 S_L(\phi)=\frac1{4\pi}\int
    \left(|d\phi|_{\widetilde g}^2+
          Q\widetilde R\phi+4\pi\mu e^{2b\phi}\right)
                         d\operatorname{vol}_{\widetilde g},
 \qquad Q=b+b^{-1}.
\]
The stress tensor and primary weights are
\[
 T(z)=-(\partial\phi)^2+Q\partial^2\phi,\qquad
 c=1+6Q^2,\qquad
 \Delta_\alpha=\alpha(Q-\alpha)
\]
for $V_\alpha=e^{2\alpha\phi}$. These formulas retain the background charge.

Set $\phi_c=2b\phi$ and hold $\lambda=\pi\mu b^2$ fixed as $b\to0$. Then the
leading rescaled action is
\[
 b^2S_L=
 \frac1{16\pi}\int
   \left(|d\phi_c|^2+2\widetilde R\phi_c+
                    16\lambda e^{\phi_c}\right)
                   d\operatorname{vol}_{\widetilde g}
   +O(b^2),
\]
with the corresponding boundary terms when the surface is represented by a
punctured or unbounded coordinate chart. In a flat coordinate its equation
of motion is
\[
 \partial\bar\partial\phi_c=2\lambda e^{\phi_c}.
\]

The exponential interaction survives at fixed nonzero $\lambda$. Thus this
semiclassical limit is not a free-boson equation. Sending $b$ to zero also
sends $c$ to infinity. It does not by itself construct a contraction of
Virasoro conformal blocks to free-boson correlators.

The distinction is visible in primary weights. For a heavy scaling
$\alpha=\eta/b$,
\[
 \Delta_\alpha=\frac{\eta(1-\eta)}{b^2}+\eta.
\]
For a light scaling $\alpha=b\sigma$,
\[
 \Delta_\alpha=\sigma+b^2\sigma(1-\sigma).
\]
An unscaled nonzero $\alpha$ has a different asymptotic behavior. A proposed
limiting operator must specify which scaling it uses.

The original Alday--Gaiotto--Tachikawa correspondence relates specified
four-dimensional gauge theories and their matter to Liouville blocks, with
the required $U(1)$ factors and parameter conventions. The relation
$b^2=\epsilon_1/\epsilon_2$ is part of such conventions. Applying it does
not turn every rank-one instanton sum into an oscillator trace. An
arithmetic version would require its own theory, comparison functor, and
controlled parameter limit. The positive geometric constructions and the
failed fixed-parameter limit leave that question open.

\section{Tate integrals and meromorphic lines}

Tate theory supplies local integrals, Euler products, and a global
functional equation. These are precise scalar constructions. We first retain
them on their actual domains, then describe what geometric sewing would
require in addition.

Normalize multiplicative Haar measure on $\mathbb Q_p^\times$ by
$\operatorname{vol}(\mathbb Z_p^\times)=1$. For an unramified character
$\chi_p$ with $\chi_p(p)=\alpha_p$, shell integration gives
\[
 Z_p(\mathbf1_{\mathbb Z_p},\chi_p,s)
  :=\int_{\mathbb Q_p^\times}
       \mathbf1_{\mathbb Z_p}(x)\chi_p(x)|x|^s\,d^\times x
  =\sum_{j\geq0}\alpha_p^jp^{-js}
  =\frac1{1-\alpha_pp^{-s}}
\]
when $|\alpha_pp^{-s}|<1$. Using $\mathbf1_{\mathbb Z_p^\times}$ instead
gives one. Thus changing the local test function can remove the unramified
Euler factor.

At the real place, with $dx/|x|$ and $\Phi_\infty(x)=e^{-\pi x^2}$,
\[
 Z_\infty(\Phi_\infty,1,s)
  =2\int_0^\infty e^{-\pi x^2}x^{s-1}\,dx
  =\pi^{-s/2}\Gamma(s/2)
       \qquad(\operatorname{Re}s>0).
\]
The substitution $u=\pi x^2$ proves the last equality. For the trivial
global character, their product is
\[
 \Lambda_\zeta(s)=\pi^{-s/2}\Gamma(s/2)\zeta(s)
       \qquad(\operatorname{Re}s>1).
\]

\begin{proposition}[Finite removal and reinsertion]
\label{v4-arith:traces:euler-removal}
For a finite prime set $S$, put
\[
 P_S(s)=\prod_{p\in S}(1-p^{-s}),\qquad
 \Phi_S(s)=P_S(s)\Lambda_\zeta(s).
\]
Using the unit-shell test function at the places in $S$ gives $\Phi_S$ in
the half-plane of absolute convergence. Reinserting their standard
unramified tests multiplies by $P_S(s)^{-1}$. The identities extend
meromorphically.
\end{proposition}
\begin{proof}
The local shell computation replaces precisely the factors indexed by $S$.
Absolute convergence allows the finite factors to be separated from the
remaining product. The meromorphic continuation of both sides preserves the
identity.
\end{proof}

For $S=\{2,3\}$, replacing the two factors by specified local parameters
gives
\[
 \Lambda_\zeta(s)
  \frac{(1-2^{-s})(1-3^{-s})}
       {(1-\alpha_2 2^{-s})(1-\alpha_3 3^{-s})}.
\]
It equals $\Lambda_\zeta(s)$ when both parameters are one. Otherwise it is
the stated partial replacement. A full character $L$-function needs its
Euler factors at every place, its conductor, and the archimedean parity
factor. Arbitrary choices at two primes do not specify a global character.

\subsection{The functional equation and its normalization}

For completeness, the scalar zeta symmetry follows from the Gaussian Poisson
identity
\[
 \theta(t)=\sum_{n\in\mathbb Z}e^{-\pi n^2t}
          =t^{-1/2}\theta(t^{-1}).
\]
For $\operatorname{Re}s>1$, termwise integration gives
\[
 \Lambda_\zeta(s)
   =\frac12\int_0^\infty(\theta(t)-1)t^{s/2}\frac{dt}{t}.
\]

Split at one and apply the theta identity to the first integral. The result
is
\begin{equation}
\label{v4-arith:traces:theta-continuation}
 \Lambda_\zeta(s)=\frac1{s(s-1)}
 +\frac12\int_1^\infty(\theta(t)-1)
       \bigl(t^{s/2}+t^{(1-s)/2}\bigr)\frac{dt}{t}.
\end{equation}
The tail decays exponentially, so the integral is entire in $s$, locally
uniformly with all derivatives. This proves meromorphic continuation, simple
poles at zero and one with residues $-1$ and $1$, and
$\Lambda_\zeta(s)=\Lambda_\zeta(1-s)$.

The entire normalization is $\xi(s)=\tfrac12s(s-1)\Lambda_\zeta(s)$.

For a primitive Dirichlet character $\chi$ of conductor $f$, let
$\chi(-1)=(-1)^a$, $a\in\{0,1\}$. The completed function and root number are
\[
 \Lambda(s,\chi)
   =\left(\frac f\pi\right)^{(s+a)/2}
         \Gamma\left(\frac{s+a}{2}\right)L(s,\chi),
 \qquad
 \varepsilon(\chi)=\frac{\tau(\chi)}{i^a\sqrt f},
\]
where $\tau(\chi)=\sum_{x\bmod f}\chi(x)e^{2\pi ix/f}$. Tate's global
functional equation, with these conventions, is
\[
 \Lambda(s,\chi)
     =\varepsilon(\chi)\Lambda(1-s,\overline\chi).
\]
Its adelic proof uses Fourier transform and Poisson summation, together with
the local character and measure data. An imprimitive character requires the
corresponding removed local factors. It cannot be inserted unchanged in the
primitive formula.

These identities construct meromorphic sections of scalar lines. For
instance, let $\mathcal M$ be the field of meromorphic functions of $s$, and
attach the line $\mathcal M e_S$ to each finite set $S$. The map deleting a
prime $p\in S$ is multiplication by $(1-p^{-s})^{-1}$. Maps for different
primes commute, so this is a functor from the finite-deletion category to
meromorphic lines. The sections $\Phi_S e_S$ are compatible with deletion.

The unmarked involution is $s\mapsto1-s$. On a marked section it has the
explicit factor
\[
 \Phi_S(s)=\frac{P_S(s)}{P_S(1-s)}\Phi_S(1-s).
\]
This is an actual comparison of the specified lines and sections. Its source
is the finite-deletion category. It is not yet a category of arithmetic
surfaces with boundary sewings.

\subsection{A scattering phase and a positive measure}

The two completed zeta normalizations also give different scattering ratios.
Put
\[
 m_\Lambda(s)=\frac{\Lambda_\zeta(2s-1)}{\Lambda_\zeta(2s)},
 \qquad
 m_\xi(s)=\frac{\xi(2s-1)}{\xi(2s)}.
\]
For the standard weight-zero Eisenstein series on $SL_2(\mathbb
Z)\backslash\mathbb H$, the constant term is $y^s+m_\Lambda(s)y^{1-s}$. The
ratio using the entire function instead satisfies
\begin{equation}
\label{v4-arith:traces:scattering-normalizations}
 m_\xi(s)=\frac{s-1}{s}m_\Lambda(s).
\end{equation}
This follows by substituting $\xi(w)=\tfrac12w(w-1)\Lambda_\zeta(w)$ and
cancelling. Thus the two scalar conventions cannot be interchanged in a
claimed normalization of an intertwining operator.

Both ratios satisfy $m(s)m(1-s)=1$ meromorphically. On $s=1/2+it$, their
modulus is one wherever they are finite and nonzero. Indeed the functional
equation and real conjugation give
\[
 \xi(2s-1)=\xi(1-2it)
          =\overline{\xi(2s)},
\]
and the same calculation applies to $\Lambda_\zeta$. The rational factor
$(s-1)/s$ also has unit modulus on this line. At the removable central
value, $m_\Lambda(1/2)=-1$ and $m_\xi(1/2)=1$. The residues $-1,1$ in
\eqref{v4-arith:traces:theta-continuation} prove the first value, and
$\xi(0)=\xi(1)=1/2$ proves the second.

Write $\Gamma_{\mathbb R}(w)=\pi^{-w/2}\Gamma(w/2)$. There is a precise
Euler product for the first ratio:
\[
 m_\Lambda(s)
   =\frac{\Gamma_{\mathbb R}(2s-1)}{\Gamma_{\mathbb R}(2s)}
      \prod_p\frac{1-p^{-2s}}{1-p^{1-2s}}
          \qquad(\operatorname{Re}s>1).
\]
Both zeta products converge absolutely on this half-plane, so taking their
quotient proves the formula.

Meromorphic continuation defines its other values. It does not assert
absolute Euler-product convergence on $\operatorname{Re}s=1/2$. Ramified
representations require their actual local $L$- and epsilon-factors,
character and measure data. They are not obtained by appending arbitrary
phases to this unramified zeta ratio.

To see what phase removal proves, work on a domain of the unitary axis where
$\xi(2s)\ne0$ and put
\[
 d(s)=|\xi(2s)|^{-2},\qquad
 P(s)=2d(s),\qquad K(s)=m_\xi(s)^{-1}d(s).
\]
These are specified scalar functions. They satisfy
\begin{equation}
\label{v4-arith:traces:phase-density}
 d(s)\leq|P(s)-K(s)|\leq3d(s),\qquad
 m_\xi(s)K(s)=d(s)=P(s)/2.
\end{equation}
The triangle and reverse-triangle inequalities prove the bounds, since
$|m_\xi|=1$ and $d>0$.

Multiplication by the phase removes that phase. It leaves the factor two in
the proposed positive coefficient. Realizing either scalar as an actual
continuous-spectrum pairing remains another question.

A measured comparison must preserve more than its diagonal density. For
example, suppose $P,K$ are positive measurable functions on a measured space
$(X,\nu)$. If $|c|^2P=K$ almost everywhere, then
\[
 f\longmapsto cf:\ L^2(X,K\,d\nu)
       \xrightarrow{\sim}L^2(X,P\,d\nu)
\]
is unitary. Integrating $|cf|^2P=|f|^2K$ proves isometry. Division by the
nonzero $c$ gives its inverse. For the phase-valued $K$ above, positivity is
not available before the form and its normalization have been specified.

A proposed arithmetic extension therefore needs a bivariate Hermitian
pairing, its domain, the measured spectral decomposition, and an intertwiner
realizing the positive form. Multiplication by $m((s_1+s_2)/2)$ alone need
not preserve the Hermitian symmetry in the two parameters. A norm identity
on individual formal spectral vectors does not provide the mixed pairings or
convergence of wave packets.

On a specified decomposition $V=V_{\mathrm{disc}}\oplus V_{\mathrm{cont}}$,
one may define the correction to be the identity on the first summand. If
both pairings make the summands orthogonal, comparisons on the two summands
combine to a comparison on $V$. Expanding the pairing of $v_d+v_c$ and
$w_d+w_c$ proves this assertion: both mixed terms vanish. Without that
orthogonality, the mixed comparisons must also be given.

A selected residual line is a particularly simple case. If the placement
carries its arithmetic vector $v_0$ to a unit vector and
$B^{\mathrm{Ar}}(v_0,\sigma v_0)=1$, then both forms on $av_0,bv_0$ equal
$a\overline b$. This calculation does not identify every residual
automorphic summand with one vacuum. Nor does discreteness of a zeta zero
set place its trace in this automorphic discrete summand. That placement
requires the actual spectral comparison.

Finally, a scalar functional equation does not transport obstruction
classes. A specified linear map $F:H\to H_{\mathrm{Ar}}$ does carry a
relation $c+c^!=0$ to $F(c)+F(c^!)=0$. Applying this fact to duality classes
requires the obstruction complexes and a comparison defined for the
corrected pairing. The identity $m(s)m(1-s)=1$ supplies no such map.

\subsection{Discriminants and residues}

Let $K$ be a number field of signature $(r_1,r_2)$ and discriminant $d_K$.
In the standard Minkowski embedding $K\hookrightarrow\mathbb
R^{r_1}\times\mathbb C^{r_2}$, use Lebesgue measure $dx\,dy$ on each complex
coordinate. Then
\begin{equation}
\label{v4-arith:traces:minkowski-volume}
 \operatorname{covol}(\mathcal O_K)
      =2^{-r_2}\sqrt{|d_K|}.
\end{equation}

To prove the factor, choose an integral basis and form its complex embedding
matrix. Its determinant has absolute value $\sqrt{|d_K|}$ by the
discriminant identity. Replacing each conjugate pair by real and imaginary
parts divides the absolute determinant by two. There are $r_2$ such pairs,
giving \eqref{v4-arith:traces:minkowski-volume}.

Thus
\[
 \log\operatorname{covol}(\mathcal O_K)
    =\frac12\log|d_K|-r_2\log2.
\]
In particular $\tfrac12\log|d_K|$ is neither an integer genus nor the
half-logarithm of this covolume.

Metrics and measure conventions matter for arithmetic Euler characteristics
as well. One can define the real discriminant expression
\[
 D_K=\log|d_K|-r_1\log(\pi/4)-r_2\log(\pi/2).
\]
It is a specified normalization of discriminant data. Identifying it with a
metrized determinant of cohomology or an arithmetic Euler characteristic
requires the actual complex, metrics and comparison formula. The expression
alone supplies none of them.

The analytic class-number formula gives the positive residue
\[
 \operatorname*{Res}_{s=1}\zeta_K(s)
  =\frac{2^{r_1}(2\pi)^{r_2}h_KR_K}
          {w_K\sqrt{|d_K|}}.
\]
Here $h_K$ is the ideal class number, $R_K$ the regulator, and $w_K$ the
number of roots of unity. The standard rank-zero regulator convention is
$R_K=1$. This is a scalar residue. For example, for $K=\mathbb Q(i)$ it is
$\pi/4$, using $h_K=R_K=1$, $w_K=4$, and $|d_K|=4$. It is therefore not in
general an integer dimension of conformal blocks.

\section{The maps required for sewing}

The finite Frobenius construction supplied an inverse pairing and explicit
gluing maps. For a family of geometric or arithmetic objects, sewing
requires their analogues, compatible under repeated gluing.

One possible coefficient for an anomaly is a determinant line. For a bounded
complex $C^\bullet$ of finite-dimensional spaces, define
\[
 \det C^\bullet
   =\bigotimes_i(\det C^i)^{(-1)^i}.
\]
Odd exponents mean dual lines. Cancellation of the boundary and coboundary
subspaces gives its canonical identification with $\bigotimes_i(\det
H^i(C))^{(-1)^i}$. Exact sequences or exact triangles give corresponding
determinant-line isomorphisms with their graded signs.

A trivialization, compatible with those isomorphisms, is additional data. It
cannot be replaced just by an orientation word or by a scalar functional
equation.

\subsection{Intermediate labels and coends}

The intermediate objects in sewing belong to a specified category. For a
small $k$-linear category $\mathcal C$, a right module $F:\mathcal
C^{\mathrm{op}}\to\operatorname{Vect}_k$ and a left module $G:\mathcal
C\to\operatorname{Vect}_k$ have a coend
\[
 F\otimes_{\mathcal C}G
    =\int^{a\in\mathcal C}F(a)\otimes_kG(a).
\]
Concretely it is the quotient of $\bigoplus_a F(a)\otimes_kG(a)$ by the
relations $F(f)(x)\otimes y=x\otimes G(f)(y)$ for $f:a\to b$, $x\in F(b)$
and $y\in G(a)$. In a dg setting, the derived coend is computed by its bar
construction with all strings of composable morphisms.

If $\mathcal C$ is split finite semisimple, the expression reduces to a
finite sum over its simples with the specified endomorphism-field balancing.
Without semisimplicity, replacing the derived coend by a sum over simple
isomorphism classes can lose extensions and higher Tor. A measured direct
integral instead requires a measurable family, a measure, Hilbert
completions and a convergence theorem. It is a different construction.

For local Weil--Deligne data, one must fix the local field, coefficient
field, rank, topology, Frobenius convention, monodromy, and allowable
inertia types. A list of all such isomorphism classes does not specify a
coend or a measure. Likewise local Langlands does not, by itself, construct
a sewing functor on arithmetic surfaces or identify its objects with WZW
labels.

The finite and archimedean actions also have different types. A change of
coefficient embedding can transport Galois labels when its semilinear maps
are given. A Hecke correspondence acts on its specified automorphic or
sheaf-theoretic carrier. The quotient $G(\mathbb A)/G(K)$ is not a group in
general, so it cannot serve as a symmetry group without a separate action
construction. Inertia and wild inertia remain part of ramified local data.

\begin{question}[Arithmetic sewing]
\label{v4-arith:traces:sewing-question}
Construct the source category of marked arithmetic objects and its gluing
operations. Give the intermediate coefficient category, the two variances of
the boundary modules, and the derived coend or measured integral that
supplies their gluing. Construct its map to the coefficient object of the
glued arithmetic space. Require compatible units, dualities, determinant
lines and repeated-sewing isomorphisms.

Compare the resulting scalar evaluations with the Tate sections on their
common domains, including the conductor, archimedean and local normalization
factors. Both the separating sewing and the self-gluing that creates a
handle are part of this question.
\end{question}

The meromorphic-line functor above provides finite Euler-factor
compatibility and functional-equation symmetry. Those results leave the
stated geometric maps to be constructed. Conversely, a geometric sewing
theorem can exist on a different specified coefficient category without yet
giving the arithmetic comparison.

\subsection{A finite analytic sewing construction}

There is a positive analytic example already constructed in
Theorem~\ref{v4-arith:asi:finite-analytic-realization}. It uses the
holomorphic lattice algebra $W=V_{E_8}$ and the finite marked coefficient
objects $(P,\lambda)$ of the arithmetic fusion category. For a compatible
junction $f:P\otimes_{R_0}Q\to S$, the field is
\[
 \mathcal J_f(u\otimes p,z)(v\otimes q)
    =Y_W(u,z)v\otimes f(p\otimes_{R_0}q).
\]
The coefficient tensor is balanced over the boundary algebra $R_0$. Its
duality maps and associators retain the original finite fusion data.

The sewing domain is explicit. Start with finitely many compact pointed
Riemann surfaces, possibly disconnected. Choose disjoint coordinate discs of
radii $r_j,r'_j$ at each paired marking and glue by $u_jv_j=t_j$, with
$0<|t_j|<r_jr'_j$. Each initial connected component retains an unsewn marked
point. Fix the initial conformal block, contragredient pairings, coordinates
and branches for fractional powers.

On this domain the descendant propagator sums converge absolutely and
locally uniformly, including simultaneous sewings that create handles. The
reason is that each coefficient contraction is finite-dimensional: the sewn
expression is a finite linear combination of ordinary $W$ conformal block
sewings. The lattice algebra is rational and $C_2$-cofinite, and its modules
here are finite sums of the vacuum module. Gui's sewing theorem applies with
the stated retained-marking hypothesis. Finite sums preserve its convergence
and sewing identities.

If an initial component has no unsewn mark, propagation of vacua first
inserts one. Its propagation isomorphism removes it afterward. The
hypothesis is not removed by silently omitting the point. This theorem
supplies actual analytic sewing for the constructed marked fusion
realization. It does not choose a globally normalized partition scalar or
the arithmetic coend in Question~\ref{v4-arith:traces:sewing-question}.

\section{Genus, rank, and coefficient data}

The genus of a smooth proper curve controls the rank of its first
cohomology. A field degree, a framing rank, and a Siegel parameter are
different inputs, even when their numerical values coincide.

For a smooth projective curve $C/K$ of genus $g$, the Jacobian has dimension
$g$ and its first cohomology has rank $2g$. At a complex embedding its Hodge
dimensions are $(g,g)$, and $\deg\omega_C=2g-2$. Arithmetic intersection and
determinant formulas additionally use a regular arithmetic model, metrics,
singular fiber contributions and analytic torsion. The topological formula
$2-2g$ is not by itself an arithmetic Euler characteristic or a Koszul level
identity.

Several possible coefficient constructions should therefore be kept
distinct.

\begin{enumerate}
\item A smooth proper curve gives its Jacobian's weight-one cohomology. An
arithmetic surface model permits metrized intersections and a determinant of
cohomology, subject to their actual hypotheses. \item A totally imaginary
field $K$ of degree $2g$ gives a Minkowski lattice in $\mathbb C^g$ after
selecting one embedding from each conjugate pair. Its quotient is a complex
torus. An abelian variety requires a positive integral Riemann form. In a CM
construction one specifies a CM type and polarization. This torus is not the
weight-zero motive whose $L$-function is $\zeta_K$. \item A rank-$2g$
ordinary crystalline representation has its coefficient field, Hodge--Tate
multiset, ordinary filtration and semilinear Frobenius. Its moduli dimension
depends on the imposed type and deformation conditions. Frobenius is not
automatically an operator with square $-1$. \item A theta construction on a
principally polarized abelian $g$-fold uses a chosen level and theta line.
Its level-$n$ section space has dimension $n^g$, whereas its first
cohomology has rank $2g$. Comparing these carriers requires an actual map.
The Steinberg relation $\{a,1-a\}=0$ in $K_2^M(K)$, for $a\ne0,1$, is a
field relation and does not supply that map. \item A surface factorization
construction needs both local operations and their interchange coherences.
Two separately named $E_1$ structures do not give an $E_2$ structure until
the hypotheses of the iterated-algebra construction hold. Double-residue and
completed duality comparisons must preserve the specified coefficients and
signs. \item A Siegel or other Shimura construction uses its actual
automorphic representation and $L$-group representation. For
$\mathrm{GSp}_{2g}$, the standard representation of the dual
$\mathrm{GSpin}_{2g+1}$ has dimension $2g+1$, while the spin representation
has dimension $2^g$. Neither is automatically the rank-$2g$ Jacobian
representation.
\end{enumerate}

These alternatives name different objects. They permit comparison problems.
They do not define six presentations of an already identified arithmetic
object. In particular an almost-complex structure belongs to a real tangent
bundle of an analytic space, whereas a crystalline Frobenius belongs to its
specified coefficient realization. An arithmetic scheme does not acquire the
former from the latter by notation.

\subsection{Weights and the zero-rank case}

For a weight-$m$ realization in motivic normalization, the expected
functional-equation involution is $s\mapsto m+1-s$. Its existence and
normalization must be established for the chosen $L$-function. The Dedekind
zeta function is weight zero and has $s\mapsto1-s$, independently of
$[K:\mathbb Q]$. The Jacobian $H^1$ function is weight one and, when its
completed functional equation is available, uses $s\mapsto2-s$. A spin or
standard automorphic function has the shifts and conductor of its own chosen
normalization.

Gamma factors likewise come from the Hodge or archimedean representation.
The number $g$ alone does not select them. For example, with
$\Gamma_{\mathbb R}(s)=\pi^{-s/2}\Gamma(s/2)$ and $\Gamma_{\mathbb
C}(s)=2(2\pi)^{-s}\Gamma(s)$, the Dedekind completion is
\[
 \Lambda_K(s)=|d_K|^{s/2}
       \Gamma_{\mathbb R}(s)^{r_1}
       \Gamma_{\mathbb C}(s)^{r_2}\zeta_K(s).
\]
A curve's weight-one completion also requires its conductor and all local
factors, including bad reduction.

At rank zero the determinant $\det(1-tF\mid0)$ is one and
$\operatorname{Sym}(0)=k$. Its all-charge trace is one. For $\mathbb P^1$,
the Jacobian is trivial and $H^1=0$. Its $H^1$ Euler factor is also one.

The full zeta function of $\mathbb P^1_{\mathbb F_q}$ is instead
\[
 Z(\mathbb P^1_{\mathbb F_q},t)
    =\frac1{(1-t)(1-qt)},
\]
because $\#\mathbb P^1(\mathbb F_{q^n})=q^n+1$. It includes degrees zero and
two. Over a number field the corresponding finite Euler product is
$\zeta_K(s)\zeta_K(s-1)$ in its initial half-plane. None of these is the
$H^1$ factor.

A number field cannot have degree zero. Thus setting $2g=0$ in a totally
imaginary degree-$2g$ family does not produce $\mathbb Q$. A Riemann Tate
sector or a rank-two Wilson label can be added as separate data. The empty
tensor contributes only its scalar unit.

For positive rank, specify the Frobenius carrier as well. If $V_\ell J$ is
the Tate module of an abelian variety, its dual $H^1_{\mathrm{\acute
et}}(J,\mathbb Q_\ell)$ has the dual action. Arithmetic and geometric
Frobenius are inverse. On the cohomological geometric convention for weight
one, polarization pairs eigenvalues with product $q$. On the Tate-module
geometric convention the product is $q^{-1}$. These distinctions persist
through the determinant trace and any Wilson comparison.

\section{Naturality and stabilization}

A genus-wise identity of scalar functions does not define maps between
coefficient objects of different rank. We state the chain-level construction
that does give such comparisons when its maps are known.

Let $\mathcal D$ be a specified category of arithmetic coefficient systems.
An object includes its local complexes, reference vectors, local actions,
completed or algebraic tensors, and all merger and residue maps being used.
Morphisms preserve these data. If conformal vectors or pairings are part of
the structure, their preservation is an additional equation. Composition of
such morphisms remains compatible because its defining squares compose. No
compatibility of a completion is inferred solely from the presence of
almost-everywhere reference vectors.

Suppose two stage constructions have differentials
$D^{(a)}_S=D^{(a)}_{0,S}+\tau^{(a)}_S$, $a=1,2$, with
\[
 (D^{(a)}_{0,S})^2=0,\qquad
 D^{(a)}_{0,S}\tau^{(a)}_S+
 \tau^{(a)}_SD^{(a)}_{0,S}
       +(\tau^{(a)}_S)^2=0.
\]
The second identity is the precise condition for the total differential to
square to zero. Its validity for two complexes does not yet give a map
between them.

\begin{proposition}[Comparison from natural stage maps]
\label{v4-arith:traces:stage-comparison}
Over a common coefficient ring, suppose quasi-isomorphisms $\Xi_{\mathsf
A,S}:C^{(1)}_S(\mathsf A) \to C^{(2)}_S(\mathsf A)$ are given. Require
naturality in $\mathsf A$, compatibility with the stage transitions, and
\[
 \Xi D^{(1)}_{0,S}=D^{(2)}_{0,S}\Xi,\qquad
 \Xi\tau^{(1)}_S=\tau^{(2)}_S\Xi.
\]
Then their algebraic filtered colimits are naturally quasi-isomorphic. Any
additional completed comparison requires preservation of this
quasi-isomorphism by the chosen completion.
\end{proposition}
\begin{proof}
Adding the two equations gives the chain-map identity. Transition
compatibility defines the colimit map. Filtered colimits of modules are
exact, so homology commutes with these colimits and the map remains a
quasi-isomorphism. The naturality squares descend from those at each stage.
The final assertion specifies the extra limiting operation and its needed
property.
\end{proof}

For ordered, arithmetic-chiral, Deninger-type and bulk-boundary
constructions, the complexes and stage maps must each be given on their
independent carriers. Defining one assignment to be another produces an
identity natural transformation, but not an independent geometric
comparison. An independently defined bulk theory also needs its boundary
functor and the comparison of its operations. Frobenius and duality pairings
are transported only when the specified maps preserve them.

\subsection{Changing the rank}

For categories $\mathcal D_g$, supply an actual stabilization functor
$S_g:\mathcal D_g\to \mathcal D_{g+1}$. Let $F_{\star,g}$ be specified chain
functors and $V_2$ a fixed graded coefficient space with zero differential.
Suppose there are natural quasi-isomorphisms
\[
 q_{\star,g}:F_{\star,g}(\mathsf A)\otimes V_2
       \longrightarrow
       F_{\star,g+1}(S_g\mathsf A).
\]
Over a field, algebraic tensoring with $V_2$ gives
\[
 H(F_{\star,g}\otimes V_2)
     \simeq H(F_{\star,g})\otimes V_2.
\]

One proof splits each complex as its cohomology and a contractible complex.
Tensor its contracting homotopy with $1_{V_2}$. For other coefficients or
completed tensors, the corresponding K\"unneth statement is an additional
hypothesis.

If the maps $q_{\star,g}$ commute with the comparison maps $\eta$ up to
specified chain homotopies, then their cohomology isomorphisms commute with
$\eta$. In particular,
\[
 \eta_{g+1}q_{\mathrm{ch},g}
   \simeq q_{\partial,g}(\eta_g\otimes1)
\]
becomes equality on homology. The same argument transports specified
Frobenius actions, retractions and pairings, provided their compatibility is
included. This gives uniformity for every $g$ where all those data exist.

An oscillator splitting alone does not construct the maps $q_{\star,g}$ for
a differential whose differential includes maps coupling residues at different places. Even a trivial action does
not imply flatness: for $\Lambda=\mathbb Z_p[[T]]$ and $M=\Lambda/(T)$, the
resolution
\[
 0\longrightarrow\Lambda
   \xrightarrow{T}\Lambda
   \longrightarrow M\longrightarrow0
\]
gives $\operatorname{Tor}_1^\Lambda(M,M)=M$. The $p$-primary Iwasawa
$\mu$-invariant of this module is zero, but $M$ is not $\Lambda$-flat. In
several variables the corresponding completion and duality hypotheses must
still be proved.

A block-diagonal group map $GL_{2g}\times GL_2 \hookrightarrow GL_{2g+2}$
also does not give a tensor decomposition of the ambient current algebra.
The latter has off-diagonal matrix currents absent from the two block Lie
algebras. A retraction of coefficient complexes requires a chosen splitting
and a chain map. It is not an automatic retraction of vertex algebras.

\subsection{Three different diagonal constructions}

For a rank-$r$ Heisenberg vertex algebra $V_r$ at level one, the tensor
$V_r\otimes V_r$ has rank $2r$ and standard central charge $2r$. Writing the
two current families as $J_i^{(1)}$ and $J_i^{(2)}$, their unscaled sums
obey
\[
 (J_i^{(1)}+J_i^{(2)})(z)
 (J_j^{(1)}+J_j^{(2)})(w)
       \sim\frac{2\delta_{ij}}{(z-w)^2}.
\]
They therefore do not give a level-one diagonal map. The normalized sums
$(J_i^{(1)}+J_i^{(2)})/\sqrt2$ do give a Heisenberg embedding, but its image
conformal vector is the Sugawara vector of only the diagonal currents. It is
different from the full tensor conformal vector: the complementary
anti-diagonal currents contribute central charge $r$.

The coordinatewise product $V_r\times V_r$ is a different vertex algebra.
Its vacuum is $(1,1)$ and its conformal vector is $(\omega,\omega)$. The map
$v\mapsto(v,v)$ preserves both and all products. Its central term is
$(r/2)(1,1)$, so the central charge remains $r$. Its weight-zero space is
two-dimensional, unlike that of the rank-$2r$ Heisenberg tensor product.
Neither construction is the free product merely by its notation.

Nor is the reduced bar an additive construction on arbitrary coefficient
algebras. Let $A=k\oplus kx$ and $B=k\oplus ky$ have square-zero
augmentation ideals. The square-zero algebra $C=k\oplus(kx\oplus ky)$ has
zero reduced-bar differential. Its bar length-two term has basis
\[
 [x|x],\quad[x|y],\quad[y|x],\quad[y|y].
\]

The direct sum of the positive reduced bars of $A$ and $B$ has only the
first and last terms in that length. The mixed terms survive in homology.
Thus an additive diagonal claim for bar constructions needs the actual
carrier and comparison. It cannot follow from linearity of one residue
operator.

\subsection{Evaluating a duality defect}

The same requirement applies to a proposed all-genus level identity. Suppose
an independently defined defect $\mathfrak d_g$ has a specified scalar
evaluation $\mathcal E_g$ such that
\[
 \mathcal E_g(\mathfrak d_g)(s)
   =\Lambda_g(s)
       -\varepsilon_g\Lambda_g(w_g-s).
\]
If the functional equation holds, the evaluated defect vanishes. Concluding
$\mathfrak d_g=0$ requires that $\mathcal E_g$ detect zero on the actual
defect space. This criterion is immediate from injectivity, and preserves
the need to construct both the evaluation and detection. Euler factors,
weights, and a scalar functional equation alone do not provide them.

\section{The level-dependent Hecke trace}

A Hecke trace on a finite space of cusp forms is another concrete scalar.
Its relation to a categorical arithmetic trace requires a map carrying the
chosen endomorphism to that Hecke operator. The classical trace itself
retains the level in both its elliptic and cusp terms.

Let $N\geq1$, $k\geq2$, and $n\geq1$. Write
\[
 \Gamma_1(N)=
 \left\{\begin{pmatrix}a&b\\c&d\end{pmatrix}
  \in SL_2(\mathbb Z):
     a\equiv d\equiv1,\ c\equiv0\pmod N\right\}.
\]
Let $S_k(\Gamma_1(N))$ be its holomorphic cusp forms and $T_n$ the usual
Hecke operator. One precise normalization is obtained by decomposing under
diamond characters $\chi$ and using on $S_k(\Gamma_0(N),\chi)$
\[
 f|T_n=
  \sum_{\sigma\in\Gamma_0(N)\backslash\Delta_n}
    (\det\sigma)^{k-1}\chi(a_\sigma)\,f|_k\sigma,
\]
where
\begin{align*}
 \Delta_n&=
 \{\sigma\in M_2(\mathbb Z):
       \det\sigma=n,\ N\mid c_\sigma,\ (a_\sigma,N)=1\},\\
 f|_k\begin{pmatrix}a&b\\c&d\end{pmatrix}(z)
       &=(cz+d)^{-k}f\left(\frac{az+b}{cz+d}\right).
\end{align*}
This convention includes indices not coprime to $N$.

Define the polynomials $P_j(t,n)$ by
\[
 \frac1{1-tx+nx^2}
      =\sum_{j\geq0}P_j(t,n)x^j.
\]
Thus $P_0=1$, $P_1=t$, and $P_j=tP_{j-1}-nP_{j-2}$. If $\alpha+\beta=t$ and
$\alpha\beta=n$, then $P_{k-2}(t,n)=
(\alpha^{k-1}-\beta^{k-1})/(\alpha-\beta)$ when the roots are distinct. The
polynomial defines the repeated-root value.

Let $\varphi$ be Euler's totient, and put
\[
 I(N)=N\prod_{p\mid N}(1+p^{-1}),\qquad
 \sigma_{1,N}(n)=
     \sum_{\substack{d\mid n\\(d,N)=1}}\frac nd.
\]
For $u\mid N$ and $u^2\mid t^2-4n$, define
\[
 B_N(u,t,n)=
  \begin{cases}
    I(N)/I(N/u),&Nu\mid t-n-1,\\
    0,&\text{otherwise},
  \end{cases}
 \qquad
 D_N(u,t,n)=\sum_{d\mid u}\mu(d)B_N(u/d,t,n).
\]
Here $\mu$ is the M\"obius function. For integers $a,d$, set
\[
 \Psi_N(a,d)=
 \sum_{\substack{rs=N\\r\mid a-1,\ s\mid d-1}}
       \varphi((r,s))\varphi(N/(r,s)).
\]

For $D>0$, let $H(D)$ be the Hurwitz class number of positive definite
integral binary quadratic forms of discriminant $-D$, including imprimitive
forms. A class has weight the reciprocal of its stabilizer order in
$PSL_2(\mathbb Z)$. Thus $H(3)=1/3$, $H(4)=1/2$. Set $H(0)=-1/12$, and set
$H(D)=0$ for a nonintegral argument or an inadmissible discriminant.

\begin{theorem}[Eichler--Selberg at level $\Gamma_1(N)$]
\label{v4-arith:traces:hecke-level}
With the preceding conventions,
\begin{align*}
 \operatorname{Tr}(T_n|S_k(\Gamma_1(N)))
  ={}&\delta_{k,2}\sigma_{1,N}(n)\\
 &-\frac{\varphi(N)}2
   \sum_{\substack{t\in\mathbb Z\\t^2\leq4n}}
     P_{k-2}(t,n)
     \sum_{\substack{u\mid N\\u^2\mid t^2-4n}}
       H\left(\frac{4n-t^2}{u^2}\right)D_N(u,t,n)\\
 &-\frac14\sum_{\substack{a,d>0\\ad=n}}
      \min(a,d)^{k-1}
       \bigl(\Psi_N(a,d)+(-1)^k\Psi_N(-a,-d)\bigr).
\end{align*}
This holds for all $n\geq1$, including bad-level indices.
\end{theorem}

The formula is Popa's Theorem 5 in \emph{On the trace formula for Hecke
operators on congruence subgroups, II}. Its elliptic level factor is $D_N$.
The terms $t^2=4n$ include the scalar contribution through $H(0)$. The last
line and the weight-two term retain the cusp and Eisenstein corrections.

For $N=1$, $D_1=1$, $\Psi_1=1$, and the even-weight formula reduces to
\[
 \delta_{k,2}\sigma_1(n)
 -\frac12\sum_{t^2\leq4n}
                P_{k-2}(t,n)H(4n-t^2)
 -\frac12\sum_{ad=n}\min(a,d)^{k-1}.
\]
At $n=1$, $k=2$, the class-number sum is $2H(0)+2H(3)+H(4)=1$. The three
contributions are $1,-1/2,-1/2$, giving $\dim S_2(SL_2(\mathbb Z))=0$. At
$k=12$, the class-number sum is $-3$, and the trace is $3/2-1/2=1$. These
finite examples check the normalization. They do not replace the
level-dependent formula.

A universal arithmetic trace family needs a coefficient space over its
actual weight and level base. Its specializations to cusp-form spaces must
commute with every $T_n$. Matching the traces does not construct the
eigenobjects, their extensions, or their categorical endomorphisms.

\section{Trace functionals and their analytic domains}

To compare an ordered bar character with an analytic trace, first turn its
homology into a scalar on a specified domain. For example, suppose $H_i$ are
Hilbert realizations and $T_i(s)$ are holomorphic trace-class families on a
half-plane $\Omega$, with $\sum_i\|T_i(s)\|_1$ locally uniformly convergent.
Then
\[
 Z(s)=\sum_i(-1)^i\operatorname{Tr}T_i(s)
\]
is holomorphic there. Continuity of trace and locally uniform convergence
prove the assertion. A regularized character needs its own prescription and
normalization. Homology itself is not this scalar.

Suppose a comparison identifies the local factors and proves $Z(s)=\xi(s)$
after the archimedean and polar normalizations. If both sides have
meromorphic continuations, the identity theorem preserves that equality. The
argument principle then computes the same zero distribution. It does not add
an independent Hilbert spectral realization.

\subsection{A common compact test class}

Take an even $f\in C_c^\infty(\mathbb R)$ and define
\[
 h(z)=\int_{\mathbb R}f(u)e^{izu}\,du,\qquad
 H(s)=h(-i(s-\tfrac12)).
\]
The function $h$ is entire and rapidly decreasing in each horizontal strip.
With $\psi_\Gamma=\Gamma'/\Gamma$, the compact Weil explicit formula is
\begin{align*}
 \sum_\rho H(\rho)={}&h(i/2)+h(-i/2)\\
 &+\frac1{2\pi}\int_{\mathbb R}h(t)
       \bigl(\operatorname{Re}\psi_\Gamma(1/4+it/2)
                              -\log\pi\bigr)\,dt\\
 &-\sum_p\sum_{j\geq1}
       (\log p)p^{-j/2}
       \bigl(f(j\log p)+f(-j\log p)\bigr).
\end{align*}

Zeros are counted with multiplicity, without assuming they are on the
critical line. The prime sum is finite because $f$ is compactly supported.
Rapid decay in strips and the classical $O(T\log(T+2))$ zero count give
absolute convergence of the zero sum.

The gamma integral converges by its logarithmic growth.

For a positively oriented rectangle $\Gamma_T$ avoiding zeros on its
boundary, the argument principle gives
\[
 \frac1{2\pi i}
    \int_{\Gamma_T}H(s)\frac{\xi'(s)}{\xi(s)}\,ds
   =\sum_{\rho\ \mathrm{inside}\ \Gamma_T}H(\rho).
\]
The entire function $\xi$ is nonzero at zero and one. Its logarithmic
derivative has no polar residues there. If one instead uses $\Lambda_\zeta$,
its poles at zero and one contribute with negative multiplicity. The choice
determines the boundary terms.

When an independent arithmetic construction has $Z=\xi$ on the proved common
domain, its argument-principle functional is the one just displayed. A
comparison with a distributional scaling trace also requires equality of
that trace on this same test class, including the gamma, prime and polar
terms. No expansion of $-\zeta'/\zeta$ on the critical line is used: its
Dirichlet series starts in $\operatorname{Re}s>1$.

\subsection{A finite trace difference with a negative value}

Even complete unitary scattering does not force a positive trace
distribution. On $L^2(\mathbb R,d\lambda)\oplus\mathbb C$, take
\[
 A_0=M_\lambda\oplus0,\qquad
 A_1=M_\lambda\oplus a,\qquad a>0,
\]
with common maximal domain $\operatorname{Dom}M_\lambda\oplus\mathbb C$.
Their difference is positive and rank one.

Both absolutely continuous subspaces are the first summand. The wave
operators are the identity there, so scattering is complete and its
scattering operator is one.

Nevertheless,
\[
 \operatorname{Tr}(h(A_1)-h(A_0))
       =h(a)-h(0)
       =\int_0^a h'(\lambda)\,d\lambda.
\]

The spectral shift is $\mathbf1_{(0,a)}$. The corresponding trace
distribution is $\delta_a-\delta_0$.

This distribution is negative on an explicit Fourier square. Choose $\phi\in
C_c^\infty(\mathbb R)$ with $\int\phi=1$, and put
$f(u)=\phi(u)+\phi(u-\pi/a)$. Then
\[
 h_f(z)=(1+e^{i\pi z/a})h_\phi(z),\qquad
 |h_f(a)|^2-|h_f(0)|^2=-4.
\]
The scattering phase is still $e^{-2\pi i\mathbf1_{(0,a)}}=1$ almost
everywhere. Thus its unit modulus does not determine the sign of the trace
functional.

A restriction to a smaller test space needs its own comparison. If a
Hermitian form is known on $\ker m$ for a moment functional $m$, and
$m(\psi)=1$, write $f=f_0+m(f)\psi$. Then
\[
 q(f,f)=q(f_0,f_0)
  +2\operatorname{Re}
       \bigl(\overline{m(f)}\,q(f_0,\psi)\bigr)
  +|m(f)|^2q(\psi,\psi).
\]
The complementary value and mixed pairing are needed to extend a positivity
statement. They cannot be discarded by a mean-zero reduction. The complete
arithmetic Weil form and its spectral comparison retain these terms.

\section{The enhanced arithmetic comparisons}

The constructions above provide several actual traces: a mixed complex, a
marked categorical closure, finite Frobenius contractions, Gibbs traces, and
classical Hecke traces. Each has its own coefficients, unit, completion and
evaluation. Their arithmetic comparison must retain those data.

The enhanced Hall question asks for a functor realizing the correspondence
operations and their collision coherences on an independently defined
arithmetic carrier. The marked-center question asks for the map preserving
the chosen center, object and action. Nonlinear norm transport requires its
maps on parameters and their multiplicative distribution laws. These are
different from constructing all-genus sewing.

An integral Frobenius trace comparison further requires the commuting square
between the chiral operations, their arithmetic sheaf realizations, the
Frobenius actions and the trace evaluations. An automorphy comparison must
identify the actual chiral functions with independently arithmetic ones,
including their Euler factors, analytic continuation, functional equation
and bounds. Neither follows from an equal scalar character.

These maps are the remaining mathematical questions. The finite
counterexamples specify what a comparison must preserve. The positive
constructions specify data available for such a comparison. The gluing maps,
their inverses where an equivalence is claimed, and their compatibility
under sewing still have to be constructed on the stated arithmetic domains.


\section*{References}

For the dg mixed complex and its Morita invariance, see Keller, \emph{On
differential graded categories}, Section 5.3, Theorem 5.2. The categorical
perfect-stack and Hecke character theorems have their separate sources and
hypotheses in Section~\ref{v4-arith:ic:traces}.

See Bost--Connes, \emph{Hecke algebras, type III factors and phase
transitions with spontaneous symmetry breaking in number theory}.
The phase transition is in Theorems 5 and 25 and Section 7. The spherical norm uses the unramified
Whittaker formula of Casselman--Shalika, \emph{The unramified principal
series of p-adic groups II: The Whittaker function}. For the global volume
normalization, see Lai, \emph{Tamagawa number of reductive algebraic
groups}, Section 1.7 and Theorem 7.1.

The holomorphic parameter bounds are Deligne, \emph{La conjecture de Weil
I}, Theorem 8.2, and Deligne--Serre, \emph{Formes modulaires de poids 1},
Theorem 4.1 and Corollary 4.2. The Tate integrals and functional equation
use Tate, \emph{Fourier analysis in number fields and Hecke's
zeta-functions}, in Cassels--Fr\"ohlich, \emph{Algebraic Number Theory}. The
compact explicit formula uses the normalization of
Theorem~\ref{v4-arith:w5-a:thm:weil-explicit}.

The fixed-point sum and hook specialization are Nekrasov--Okounkov,
\emph{Seiberg--Witten theory and random partitions}, arXiv:hep-th/0306238,
Sections 3.1--3.2. The matter and rank-two gauge/Liouville conventions are
Alday--Gaiotto--Tachikawa, \emph{Liouville correlation functions from
four-dimensional gauge theories}, arXiv:0906.3219, including Appendix B.1.
The geometric transfer operators and divisor theorem are Maulik--Okounkov,
\emph{Quantum groups and quantum cohomology}, arXiv:1211.1287, Theorem
1.3.1. For the Liouville action and heavy/light limits, see
Harlow--Maltz--Witten, \emph{Analytic continuation of Liouville theory},
arXiv:1108.4417, Section 2.

The analytic sewing input is Gui, \emph{Convergence of sewing conformal
blocks}, arXiv:2011.07450, Theorem 13.1, with its component-marking
hypothesis. The finite fusion and conformal realizations used here are
constructed in Chapters~\ref{v4-arith:mf:chapter}
and~\ref{v4-arith:asi:chapter}. The level-dependent Hecke formula is Popa,
\emph{On the trace formula for Hecke operators on congruence subgroups, II},
arXiv:1706.02691, Theorem 5. Its conventions include indices divisible by
the level.
